\documentclass[aspectratio=169,11pt]{beamer}

\usetheme{Madrid}
\usecolortheme{dolphin}
\setbeamertemplate{navigation symbols}{}
\setbeamertemplate{footline}{%
  \leavevmode%
  \hbox{%
    \begin{beamercolorbox}[wd=.78\paperwidth,ht=2.7ex,dp=1.2ex,leftskip=1.2em]{author in head/foot}%
      \usebeamerfont{author in head/foot}A Journey to the World of Numbers
    \end{beamercolorbox}%
    \begin{beamercolorbox}[wd=.22\paperwidth,ht=2.7ex,dp=1.2ex,rightskip=1.2em plus1fil]{date in head/foot}%
      \usebeamerfont{date in head/foot}\hfill\insertframenumber/\inserttotalframenumber
    \end{beamercolorbox}%
  }%
  \vskip0pt%
}

\definecolor{journeyblue}{RGB}{27,75,125}
\definecolor{journeygold}{RGB}{184,126,23}
\setbeamercolor{structure}{fg=journeyblue}
\setbeamercolor{alerted text}{fg=journeygold}
\setbeamercolor{title}{fg=white,bg=journeyblue}
\setbeamercolor{frametitle}{fg=white,bg=journeyblue}
\setbeamercolor{block title}{fg=white,bg=journeyblue}
\setbeamercolor{block body}{bg=journeyblue!5}
\setbeamercolor{block title alerted}{fg=white,bg=journeygold}
\setbeamercolor{block body alerted}{bg=journeygold!8}

\usepackage[T1]{fontenc}
\usepackage{lmodern}
\usepackage{amsmath,amssymb,mathtools}
\usepackage{microtype}

\newcommand{\N}{\mathbb N}
\newcommand{\suc}[1]{S(#1)}
\newcommand{\Z}{\mathbb Z}
\newcommand{\Q}{\mathbb Q}
\newcommand{\R}{\mathbb R}
\newcommand{\cl}[2]{\left[(#1,#2)\right]}
\newcommand{\qcl}[2]{\left[\frac{#1}{#2}\right]}
\newcommand{\Adic}[1]{\operatorname{Adic}(#1)}
\newcommand{\Zmod}[1]{\mathbb Z/#1\mathbb Z}

\title{A Journey to the World of Numbers}
\author{Riccardo Brasca \and Shanwen Wang}
\institute{Remnim University}
\date{July 2026}

\AtBeginSection[]{
  \begin{frame}[plain]
    \centering
    \vfill
    {\usebeamerfont{title}\color{journeyblue}\Huge\insertsectionhead}
    \vfill
  \end{frame}
}

\begin{document}

\begin{frame}[plain]
  \titlepage
\end{frame}

\begin{frame}
  \small
  \tableofcontents[hideallsubsections]
\end{frame}

\section{The natural numbers}

\begin{frame}{Starting from \(0\)}
  The natural numbers form the sequence
  \[
    0,\qquad \suc 0,\qquad \suc{\suc 0},\qquad
    \suc{\suc{\suc 0}},\quad\ldots
  \]
  governed by the following principles.

  \begin{enumerate}
    \item \(0\in\N\).
    \item If \(n\in\N\), then its \alert{successor} \(\suc n\in\N\).
    \item Every natural number is either \(0\) or the successor of a natural number.
    \item No successor equals \(0\): \(\suc n\ne0\).
    \item Successors are unique: \(\suc m=\suc n\) implies \(m=n\).
  \end{enumerate}

  We set \(1=\suc 0\).  Once addition is defined, we will prove that
  \(\suc n=n+1\).
\end{frame}

\begin{frame}{The induction principle}
  \begin{block}{Principle of mathematical induction}
    Let \(P(n)\) be a statement about \(n\in\N\).  If
    \begin{enumerate}
      \item \(P(0)\) is true, and
      \item for every \(n\in\N\), \(P(n)\) implies \(P(\suc n)\),
    \end{enumerate}
    then \(P(n)\) is true for every \(n\in\N\).
  \end{block}

  \vspace{0.4em}
  Induction mirrors the way the natural numbers are generated:
  \[
    \underbrace{P(0)}_{\text{start}}
    \Longrightarrow P(\suc 0)
    \Longrightarrow P(\suc{\suc 0})
    \Longrightarrow \cdots
  \]
\end{frame}

\begin{frame}{The recursion principle}
  \begin{block}{Defining a function on \(\N\)}
    Let \(X\) be a set, choose \(x_0\in X\), and let
    \(G\colon\N\times X\to X\) be a rule.  There is one and only one function
    \(f\colon\N\to X\) satisfying
    \[
      f(0)=x_0,
      \qquad
      f(\suc n)=G\bigl(n,f(n)\bigr)
      \quad\text{for every }n\in\N.
    \]
  \end{block}

  \begin{alertblock}{Why this matters}
    Addition, multiplication, and the predecessor function will all be defined
    by this principle.
  \end{alertblock}
\end{frame}

\subsection{Addition}

\begin{frame}{Definition of addition}
  Fix \(a\in\N\).  We define \(a+b\) by recursion on \(b\):
  \[
    \boxed{a+0=a},
    \qquad
    \boxed{a+\suc b=\suc{a+b}}.
  \]

  Thus the second input tells us how many times to take the successor of \(a\).
  For example,
  \begin{align*}
    a+\suc{\suc 0}
      &=\suc{a+\suc 0}\\
      &=\suc{\suc{a+0}}\\
      &=\suc{\suc a}.
  \end{align*}

  Since \(1=\suc 0\), the defining equations immediately give
  \[
    a+1=a+\suc 0=\suc{a+0}=\suc a.
  \]
\end{frame}

\begin{frame}{Zero on the left}
  \begin{block}{Proposition}
    For every \(a\in\N\), \(0+a=a\).
  \end{block}

  \begin{proof}
    We use induction on \(a\).

    \emph{Base case.} By the first defining equation,
    \(0+0=0\).

    \emph{Inductive step.} Suppose that \(0+n=n\).  Then
    \[
      0+\suc n
        =\suc{0+n}
        =\suc n.
    \]
    The first equality is the second defining equation for addition; the
    second uses the induction hypothesis.  Hence the statement passes from
    \(n\) to \(\suc n\).

    By induction, \(0+a=a\) for all \(a\in\N\).
  \end{proof}
\end{frame}

\begin{frame}{Moving a successor through a sum}
  \begin{block}{Proposition}
    For all \(a,b\in\N\),
    \[
      \suc a+b=\suc{a+b}.
    \]
  \end{block}

  \begin{proof}
    Fix \(a\) and use induction on \(b\).

    \emph{Base case.}
    \(\suc a+0=\suc a=\suc{a+0}\).

    \emph{Inductive step.} Assume \(\suc a+n=\suc{a+n}\).  Then
    \begin{align*}
      \suc a+\suc n
        &=\suc{\suc a+n} &&\text{by the definition of addition}\\
        &=\suc{\suc{a+n}} &&\text{by the induction hypothesis}\\
        &=\suc{a+\suc n} &&\text{by the definition of addition}.
    \end{align*}
    This proves the required identity for \(\suc n\).
  \end{proof}
\end{frame}

\begin{frame}{Associativity of addition}
  \footnotesize
  \begin{block}{Theorem}
    For all \(a,b,c\in\N\),
    \[
      (a+b)+c=a+(b+c).
    \]
  \end{block}

  \begin{proof}
    Fix \(a,b\) and use induction on \(c\).

    \emph{Base case.}
    \((a+b)+0=a+b=a+(b+0)\).

    \emph{Inductive step.} Suppose \((a+b)+n=a+(b+n)\).  Then
    \begin{align*}
      (a+b)+\suc n
        &=\suc{(a+b)+n}\\
        &=\suc{a+(b+n)}\\
        &=a+\suc{b+n}\\
        &=a+(b+\suc n).
    \end{align*}
    Here the second equality is the induction hypothesis, and the other
    equalities are the defining equation for addition, read in either direction.
  \end{proof}
\end{frame}

\begin{frame}{Commutativity of addition}
  \begin{block}{Theorem}
    For all \(a,b\in\N\), \(a+b=b+a\).
  \end{block}

  \begin{proof}
    Fix \(b\) and use induction on \(a\).

    \emph{Base case.} We already proved \(0+b=b\), while \(b+0=b\) by
    definition.  Hence \(0+b=b+0\).

    \emph{Inductive step.} Suppose \(n+b=b+n\).  Using the formula for
    moving a successor through a sum,
    \begin{align*}
      \suc n+b
        &=\suc{n+b}\\
        &=\suc{b+n} &&\text{by the induction hypothesis}\\
        &=b+\suc n &&\text{by the definition of addition}.
    \end{align*}
    Therefore the equality holds for \(\suc n\), and induction completes the proof.
  \end{proof}
\end{frame}

\begin{frame}{When can a sum be zero?}
  \small
  \begin{block}{Proposition}
    If \(a+b=0\), then \(a=0\) and \(b=0\).
  \end{block}

  \begin{proof}
    First inspect \(b\).

    If \(b=0\), then \(a+b=a+0=a\), so the assumption gives \(a=0\).

    If \(b=\suc n\), then
    \[
      a+b=a+\suc n=\suc{a+n},
    \]
    which cannot equal \(0\).  Thus the second possibility is ruled out, and
    \(b=0\); the preceding paragraph then gives \(a=0\).

    Conversely, \(0+0=0\), so in fact
    \[
      a+b=0\quad\Longleftrightarrow\quad a=0\text{ and }b=0.
    \]
  \end{proof}
\end{frame}

\begin{frame}{Cancellation for addition}
  \begin{block}{Theorem}
    If \(a+b=a+c\), then \(b=c\).
  \end{block}

  \begin{proof}
    We use induction on \(a\), proving the assertion simultaneously for all
    \(b,c\in\N\).

    \emph{Base case.} If \(0+b=0+c\), then \(b=c\), because \(0+b=b\) and
    \(0+c=c\).

    \emph{Inductive step.} Assume cancellation is valid with \(n\) on the left.
    If
    \[
      \suc n+b=\suc n+c,
    \]
    moving the successors through the sums gives
    \(\suc{n+b}=\suc{n+c}\).  Uniqueness of successors yields \(n+b=n+c\),
    and the induction hypothesis now yields \(b=c\).
  \end{proof}
\end{frame}

\subsection{Multiplication}

\begin{frame}{Definition of multiplication}
  Fix \(a\in\N\).  We define \(ab\) by recursion on \(b\):
  \[
    \boxed{a\cdot0=0},
    \qquad
    \boxed{a\cdot\suc b=ab+a}.
  \]

  Multiplication is therefore repeated addition.  For example,
  \begin{align*}
    a\cdot\suc{\suc 0}
      &=a\cdot\suc 0+a\\
      &=(a\cdot0+a)+a\\
      &=(0+a)+a\\
      &=a+a.
  \end{align*}
\end{frame}

\begin{frame}{Moving a successor through a product}
  \footnotesize
  \begin{block}{Proposition}
    For all \(a,b\in\N\),
    \[
      \suc a\cdot b=ab+b.
    \]
  \end{block}

  \begin{proof}
    Fix \(a\) and use induction on \(b\).

    \emph{Base case.}
    \(\suc a\cdot0=0=a\cdot0+0\).

    \emph{Inductive step.} Suppose \(\suc a\cdot n=an+n\).  Then
    \begin{align*}
      \suc a\cdot\suc n
        &=\suc a\cdot n+\suc a\\
        &=(an+n)+\suc a\\
        &=\suc{(an+n)+a}\\
        &=\suc{(an+a)+n}\\
        &=a\cdot\suc n+\suc n.
    \end{align*}
    Associativity and commutativity of addition justify the fourth equality.
  \end{proof}
\end{frame}

\begin{frame}{Zero on the left of a product}
  \begin{block}{Proposition}
    For every \(a\in\N\), \(0\cdot a=0\).
  \end{block}

  \begin{proof}
    We use induction on \(a\).

    \emph{Base case.} \(0\cdot0=0\) by definition.

    \emph{Inductive step.} Suppose \(0\cdot n=0\).  Then
    \[
      0\cdot\suc n=0\cdot n+0=0+0=0.
    \]
    Thus the property passes to the successor, and induction proves it for
    every \(a\).
  \end{proof}

  Together with the defining identity \(a\cdot0=0\), this says that \(0\) is
  absorbing on both sides of multiplication.
\end{frame}

\begin{frame}{Left distributivity}
  \footnotesize
  \begin{block}{Theorem}
    For all \(a,b,c\in\N\),
    \[
      a(b+c)=ab+ac.
    \]
  \end{block}

  \begin{proof}
    Fix \(a,b\) and use induction on \(c\).

    \emph{Base case.}
    \(a(b+0)=ab=ab+0=ab+a\cdot0\).

    \emph{Inductive step.} Suppose \(a(b+n)=ab+an\).  Then
    \begin{align*}
      a(b+\suc n)
        &=a\suc{(b+n)}\\
        &=a(b+n)+a\\
        &=(ab+an)+a\\
        &=ab+(an+a)\\
        &=ab+a\suc n.
    \end{align*}
    We used the recursion equations and associativity of addition.
  \end{proof}
\end{frame}

\begin{frame}{The multiplicative identity}
  \begin{block}{Proposition}
    For every \(a\in\N\), \(a\cdot1=a\).
  \end{block}

  Since \(1=\suc 0\), the definition of multiplication gives
  \[
    a\cdot1=a\cdot\suc 0=a\cdot0+a=0+a=a.
  \]

  \begin{block}{A useful observation}
    The equality \(1\cdot a=a\) will follow once multiplication is known to be
    commutative.  It can also be proved directly by induction on \(a\).

  The base case is \(1\cdot0=0\).  If \(1\cdot n=n\), then
  \[
    1\cdot\suc n=1\cdot n+1=n+1=\suc n.
  \]
  \end{block}
\end{frame}

\begin{frame}{Commutativity of multiplication}
  \begin{block}{Theorem}
    For all \(a,b\in\N\), \(ab=ba\).
  \end{block}

  \begin{proof}
    Fix \(a\) and use induction on \(b\).

    \emph{Base case.}
    \(a\cdot0=0=0\cdot a\).

    \emph{Inductive step.} Suppose \(an=na\).  Then
    \begin{align*}
      a\cdot\suc n
        &=an+a &&\text{by the definition of multiplication}\\
        &=na+a &&\text{by the induction hypothesis}\\
        &=\suc n\cdot a &&\text{by the successor--product formula}.
    \end{align*}
    Hence \(a\cdot b=b\cdot a\) for all \(b\).
  \end{proof}
\end{frame}

\begin{frame}{Right distributivity}
  \begin{block}{Theorem}
    For all \(a,b,c\in\N\),
    \[
      (a+b)c=ac+bc.
    \]
  \end{block}

  \begin{proof}
    We convert right multiplication into left multiplication, use left
    distributivity, and convert back:
    \begin{align*}
      (a+b)c
        &=c(a+b) &&\text{commutativity of multiplication}\\
        &=ca+cb &&\text{left distributivity}\\
        &=ac+bc &&\text{commutativity of multiplication}.
    \end{align*}
  \end{proof}

  This short argument illustrates an important strategy: a symmetry theorem can
  turn a new statement into one already proved.
\end{frame}

\begin{frame}{Associativity of multiplication}
  \small
  \begin{block}{Theorem}
    For all \(a,b,c\in\N\),
    \[
      (ab)c=a(bc).
    \]
  \end{block}

  \begin{proof}
    Fix \(a,b\) and use induction on \(c\).

    \emph{Base case.}
    \((ab)0=0=a0=a(b0)\).

    \emph{Inductive step.} Suppose \((ab)n=a(bn)\).  Then
    \begin{align*}
      (ab)\suc n
        &=(ab)n+ab\\
        &=a(bn)+ab\\
        &=a(bn+b)\\
        &=a(b\suc n).
    \end{align*}
    The third equality uses left distributivity in reverse.
  \end{proof}
\end{frame}

\begin{frame}{No zero divisors}
  \small
  \begin{block}{Theorem}
    If \(a\ne0\) and \(b\ne0\), then \(ab\ne0\).
  \end{block}

  Since every nonzero natural number is a successor, write
  \(a=\suc p=p+1\) and \(b=\suc q=q+1\).  Then
  \begin{align*}
    ab
      &=(p+1)(q+1)\\
      &=p(q+1)+1(q+1)\\
      &=(pq+p)+(q+1)\\
      &=(pq+p+q)+1\\
      &=\suc{pq+p+q}.
  \end{align*}
  The final expression is a successor, so it cannot be \(0\).  Therefore
  \(ab\ne0\).

  \begin{alertblock}{Equivalent form}
    If \(ab=0\), then \(a=0\) or \(b=0\).
  \end{alertblock}
\end{frame}

\begin{frame}{The predecessor function}
  Define \(\operatorname{pred}\colon\N\to\N\) by
  \[
    \operatorname{pred}(0)=0,
    \qquad
    \operatorname{pred}(\suc a)=a.
  \]

  Consequently,
  \[
    \operatorname{pred}(a+1)=a
  \]
  for every \(a\), because \(a+1=\suc a\).

  Conversely, if \(a\ne0\), then \(a\) is a successor, say \(a=\suc n\), and
  \[
    \operatorname{pred}(a)+1
      =\operatorname{pred}(\suc n)+1
      =n+1
      =\suc n
      =a.
  \]

  The condition \(a\ne0\) is essential: our definition gives
  \(\operatorname{pred}(0)+1=1\ne0\).
\end{frame}

\begin{frame}{The algebra obtained so far}
  For all \(a,b,c\in\N\), we have established
  \begin{align*}
    a+b&=b+a, & (a+b)+c&=a+(b+c), & a+0&=a,\\
    ab&=ba, & (ab)c&=a(bc), & a\cdot1&=a,\\
    a(b+c)&=ab+ac, & (a+b)c&=ac+bc, & a\cdot0&=0.
  \end{align*}

  Thus \(\N\), together with \(0,1,+,\cdot\), is a
  \alert{commutative semiring}.  In addition,
  \begin{itemize}
    \item \(0\ne1\);
    \item addition admits cancellation;
    \item a product of two nonzero numbers is nonzero.
  \end{itemize}

  We now build the order directly from addition.
\end{frame}

\subsection{The order relation}

\begin{frame}{Definition of \(\leq\)}
  \begin{block}{Definition}
    For \(a,b\in\N\), we write
    \[
      \boxed{a\le b}
      \qquad\text{if and only if}\qquad
      \boxed{\text{there is an }x\in\N\text{ such that }b=a+x.}
    \]
  \end{block}

  The number \(x\) is the \alert{gap} from \(a\) to \(b\).

  Immediate examples are
  \begin{itemize}
    \item \(0\le a\), because \(a=0+a\);
    \item \(a\le a\), because \(a=a+0\);
    \item \(a\le\suc a\), because \(\suc a=a+1\);
    \item \(a\le a+b\), because \(a+b=a+b\).
  \end{itemize}

  We must prove that this relation really has all the expected order properties.
\end{frame}

\begin{frame}{Reflexivity of \(\leq\)}
  \begin{block}{Proposition}
    For every \(a\in\N\),
    \[
      a\le a.
    \]
  \end{block}

  \begin{proof}
    By definition, we must find an \(x\in\N\) such that
    \[
      a=a+x.
    \]
    Choose \(x=0\).  The additive identity gives
    \[
      a=a+0,
    \]
    so \(0\) is a gap from \(a\) to itself.  Hence \(a\le a\).
  \end{proof}
\end{frame}

\begin{frame}{Transitivity of \(\leq\)}
  \begin{block}{Proposition}
    If \(a\le b\) and \(b\le c\), then \(a\le c\).
  \end{block}

  \begin{proof}
    Since \(a\le b\), choose \(x\in\N\) such that
    \[
      b=a+x.
    \]
    Since \(b\le c\), choose \(y\in\N\) such that
    \[
      c=b+y.
    \]
    Substituting the first equality into the second and using associativity,
    \[
      c=(a+x)+y=a+(x+y).
    \]
    Because \(x+y\in\N\), it is a gap from \(a\) to \(c\).  Hence \(a\le c\).
  \end{proof}
\end{frame}

\begin{frame}{Antisymmetry of \(\leq\): setting up the gaps}
  \begin{block}{Theorem}
    If \(a\le b\) and \(b\le a\), then \(a=b\).
  \end{block}

  Choose gaps \(x,y\in\N\) such that
  \[
    b=a+x,
    \qquad
    a=b+y.
  \]
  Substitution gives
  \[
    a=b+y=(a+x)+y=a+(x+y).
  \]
  But also \(a=a+0\).  Therefore
  \[
    a+0=a+(x+y).
  \]
  Cancellation for addition yields
  \[
    0=x+y.
  \]
\end{frame}

\begin{frame}{Antisymmetry of \(\leq\): finishing the proof}
  From \(x+y=0\), the zero-sum proposition gives \(x=0\) and \(y=0\).
  In particular, the first gap vanishes, and therefore
  \[
    b=a+x=a+0=a.
  \]
  Hence \(a=b\), as required.

\bigskip

  Notice the chain of earlier results used here:
  \[
    \text{associativity}
    \;\Longrightarrow\;
    \text{addition cancellation}
    \;\Longrightarrow\;
    \text{antisymmetry}.
  \]
\end{frame}

\begin{frame}{Totality of \(\leq\): the induction}
  \begin{block}{Theorem}
    For every \(a,b\in\N\), either \(a\le b\) or \(b\le a\).
  \end{block}

  Fix \(b\) and use induction on \(a\).

  \emph{Base case.} For \(a=0\), we have \(0\le b\), witnessed by the gap \(b\).

  \emph{Inductive step.} Suppose that either \(n\le b\) or \(b\le n\).
  We must compare \(\suc n\) with \(b\).

  If \(b\le n\), then \(b\le n\le\suc n\), so transitivity gives
  \(b\le\suc n\).

  It remains to study the case \(n\le b\).
\end{frame}

\begin{frame}{Totality of \(\leq\): analysing the remaining gap}
  Suppose \(n\le b\).  Choose \(x\in\N\) with \(b=n+x\).
  There are two possibilities.

  \begin{itemize}
    \item If \(x=0\), then \(b=n\), and hence \(b=n\le\suc n\).
    \item If \(x=\suc y\), then
      \begin{align*}
        b=n+\suc y
          &=\suc{n+y}\\
          &=\suc n+y.
      \end{align*}
      Thus \(y\) is a gap from \(\suc n\) to \(b\), so \(\suc n\le b\).
  \end{itemize}

  In every case, \(\suc n\le b\) or \(b\le\suc n\).  The inductive step is
  complete, and therefore the order is total.
\end{frame}

\begin{frame}{A useful successor test}
  \begin{block}{Proposition}
    For all \(a,b\in\N\),
    \[
      a\le\suc b
      \quad\Longleftrightarrow\quad
      a=\suc b\ \text{ or }\ a\le b.
    \]
  \end{block}

  Suppose first that \(\suc b=a+x\).
  \begin{itemize}
    \item If \(x=0\), then \(\suc b=a\), so \(a=\suc b\).
    \item If \(x=\suc y\), then
      \(\suc b=a+\suc y=\suc{a+y}\).  Uniqueness of successors gives
      \(b=a+y\), hence \(a\le b\).
  \end{itemize}

  Conversely, equality gives \(a\le\suc b\) by reflexivity; and if \(a\le b\),
  then \(a\le b\le\suc b\) by transitivity.
\end{frame}

\begin{frame}{We have a linear order}
  The relation \(\leq\) is
  \begin{itemize}
    \item \alert{reflexive}: \(a\le a\);
    \item \alert{transitive}: \(a\le b\le c\) implies \(a\le c\);
    \item \alert{antisymmetric}: \(a\le b\) and \(b\le a\) imply \(a=b\);
    \item \alert{total}: either \(a\le b\) or \(b\le a\).
  \end{itemize}

  Therefore \((\N,\le)\) is a \alert{linearly ordered set}.

  The gap definition also gives, without further work,
  \[
    a\le a+b
    \qquad\text{and}\qquad
    a\le b+a,
  \]
  because addition is commutative.
\end{frame}

\subsection{Strict order and arithmetic}

\begin{frame}{Definition of strict inequality}
  \begin{block}{Definition}
    For \(a,b\in\N\),
    \[
      a<b
      \quad\Longleftrightarrow\quad
      a\le b\ \text{ and }\ b\nleq a.
    \]
  \end{block}

  \begin{block}{Equivalent form}

    \[
      a<b
      \quad\Longleftrightarrow\quad
      a\le b\ \text{ and }\ a\ne b.
    \]
  \end{block}

  Indeed, if \(a<b\) and \(a=b\), reflexivity would give \(b\le a\), a
  contradiction.  Conversely, if \(a\le b\), \(a\ne b\), and \(b\le a\),
  antisymmetry would force \(a=b\), again a contradiction.
\end{frame}

\begin{frame}{Positive numbers}
  \begin{block}{Proposition}
    For every \(a\in\N\),
    \[
      a\ne0\quad\Longleftrightarrow\quad 0<a.
    \]
  \end{block}

  \begin{proof}
    We always have \(0\le a\).  If \(a\ne0\), the equivalent form of strict
    inequality therefore gives \(0<a\).

    Conversely, if \(0<a\), then strict inequality implies \(0\ne a\), which
    is the same as \(a\ne0\).
  \end{proof}

  Thus, for natural numbers, \alert{nonzero} and \alert{positive} mean exactly
  the same thing.
\end{frame}

\begin{frame}{Strict inequality means a positive gap}
  \begin{block}{Theorem}
    For \(a,b\in\N\),
    \[
      a<b
      \quad\Longleftrightarrow\quad
      \text{there is an }x\ne0\text{ such that }b=a+x.
    \]
  \end{block}

  If \(a<b\), choose \(x\) with \(b=a+x\).  If \(x=0\), then
  \(b=a+0=a\), contradicting \(a\ne b\).  Hence \(x\ne0\).

  Conversely, suppose \(b=a+x\) with \(x\ne0\).  Then \(a\le b\).  If \(a=b\),
  we would have
  \[
    a+0=a=b=a+x.
  \]
  Cancellation would give \(x=0\), a contradiction.  Thus \(a\ne b\), and
  consequently \(a<b\).
\end{frame}

\begin{frame}{Addition preserves order}
  \begin{block}{Theorem}
    If \(a\le b\), then \(a+c\le b+c\).
  \end{block}

  Choose \(x\in\N\) with \(b=a+x\).  We look for a gap from \(a+c\) to \(b+c\).
  Using associativity and commutativity,
  \begin{align*}
    b+c
      &=(a+x)+c\\
      &=a+(x+c)\\
      &=a+(c+x)\\
      &=(a+c)+x.
  \end{align*}
  Thus the same \(x\) is a gap from \(a+c\) to \(b+c\), proving
  \(a+c\le b+c\).

  By commutativity, adding \(c\) on the left also preserves order.
\end{frame}

\begin{frame}{Cancelling the same summand from an inequality}
  \begin{block}{Theorem}
    If \(a+b\le a+c\), then \(b\le c\).
  \end{block}

  Choose \(x\in\N\) such that
  \[
    a+c=(a+b)+x.
  \]
  By associativity, the right-hand side is \(a+(b+x)\), so
  \[
    a+c=a+(b+x).
  \]
  Addition cancellation gives
  \[
    c=b+x.
  \]
  Hence \(x\) is a gap from \(b\) to \(c\), and \(b\le c\).

  Together with order preservation, this yields
  \[
    a+b\le a+c\quad\Longleftrightarrow\quad b\le c.
  \]
\end{frame}

\begin{frame}{Positive multiplication preserves strict order}
  \begin{block}{Theorem}
    If \(0<c\) and \(a<b\), then \(ca<cb\).
  \end{block}

  Because \(a<b\), write \(b=a+x\) with \(x\ne0\).  Because \(0<c\), we also
  have \(c\ne0\).  Distributivity gives
  \[
    cb=c(a+x)=ca+cx.
  \]
  Since \(c\ne0\) and \(x\ne0\), the no-zero-divisors theorem gives
  \(cx\ne0\).  Thus \(cb\) is obtained from \(ca\) by adding the nonzero gap
  \(cx\).  The positive-gap characterization now gives
  \[
    ca<cb.
  \]
\end{frame}

\begin{frame}{Multiplication on the right}
  \begin{block}{Corollary}
    If \(0<c\) and \(a<b\), then \(ac<bc\).
  \end{block}

  By commutativity of multiplication,
  \[
    ac=ca<cb=bc.
  \]

  \begin{alertblock}{Why positivity is necessary}
    If \(c=0\), then \(ac=0=bc\), even when \(a<b\).  Multiplication by zero
    destroys strict inequality.
  \end{alertblock}

  For weak inequalities, no positivity assumption is needed: if \(a\le b\),
  write \(b=a+x\); then \(bc=ac+xc\), so \(ac\le bc\).
\end{frame}

\begin{frame}{Cancellation for multiplication: reduction to order}
  \begin{block}{Theorem}
    Suppose \(a\ne0\).  If \(ab=ac\), then \(b=c\).
  \end{block}

  Since the order is total, either \(b\le c\) or \(c\le b\).

  Assume first that \(b\le c\).  If \(b\ne c\), then \(b<c\).  Since
  \(a\ne0\) is equivalent to \(0<a\), strict order preservation gives
  \[
    ab<ac,
  \]
  contradicting \(ab=ac\).  Therefore \(b=c\).

  The case \(c\le b\) is symmetric: if \(c\ne b\), then \(c<b\), hence
  \(ac<ab\), again contradicting the assumed equality.
\end{frame}

\begin{frame}{Right cancellation for multiplication}
  \begin{block}{Corollary}
    Suppose \(a\ne0\).  If \(ba=ca\), then \(b=c\).
  \end{block}

  Commutativity converts the hypothesis into
  \[
    ab=ac.
  \]
  The multiplication cancellation theorem therefore gives \(b=c\).

  \begin{alertblock}{Why \(a\ne0\) is necessary}
    For \(a=0\), the equality \(b0=c0\) holds for every \(b,c\in\N\), so no
    conclusion about \(b\) and \(c\) is possible.
  \end{alertblock}
\end{frame}

\subsection{What we have built}

\begin{frame}{Final picture}
  Starting from the natural-number principles, we proved that
  \begin{itemize}
    \item addition and multiplication obey the usual associative,
      commutative, distributive, and identity laws;
    \item sums and products have strong cancellation properties;
    \item nonzero products are nonzero;
    \item the relation \(a\le b\iff b=a+x\) is a linear order;
    \item \(a<b\) means exactly that the gap from \(a\) to \(b\) is positive;
    \item addition preserves order, and multiplication by a positive number
      preserves strict order.
  \end{itemize}

\end{frame}

\begin{frame}{An equation without solutions in \(\N\)}
  Suppose that some \(x\in\N\) satisfied
  \[
    3+x=1.
  \]
  By commutativity and the recursive definition of addition,
  \[
    3+x=x+3=\suc{\suc{\suc x}},
    \qquad
    1=\suc 0.
  \]
  Hence
  \[
    \suc{\suc{\suc x}}=\suc 0.
  \]
  Injectivity of the successor operation would then give
  \[
    \suc{\suc x}=0,
  \]
  which is impossible because zero is not a successor.  Therefore
  \(3+x=1\) has no solution in \(\N\).
\end{frame}

\section{The integers}

\subsection{Formal differences}

\begin{frame}{Why introduce the integers?}
  In \(\N\), the equation
  \[
    3+x=1
  \]
  has no solution.  We want a larger number system in which it has a solution
  (namely \(x=-2\)).

  Every integer can be thought of as a difference
  \[
    a-b,\qquad a,b\in\N.
  \]
  Since subtraction is not yet available, we begin with the ordered pair
  \((a,b)\), intended to represent that difference.

  \begin{alertblock}{The central difficulty}
    The same integer has many representations:
    \[
      5-2=4-1=3-0.
    \]
    We must decide exactly when two pairs represent the same difference.
  \end{alertblock}
\end{frame}

\begin{frame}{Equality of formal differences}
  \begin{block}{Definition}
    For \(a,b,c,d\in\N\), declare
    \[
      (a,b)\sim(c,d)
      \quad\Longleftrightarrow\quad
      a+d=b+c.
    \]
  \end{block}

  The equation uses only addition in \(\N\), but it expresses
  \[
    a-b=c-d.
  \]
  Indeed, cross-adding \(b+d\) to the informal equation gives precisely
  \(a+d=b+c\).

  For example,
  \[
    (5,2)\sim(4,1)
    \qquad\text{because}\qquad
    5+1=2+4.
  \]
\end{frame}

\begin{frame}{Reflexivity and symmetry}
  \begin{block}{Proposition}
    The relation \(\sim\) is reflexive and symmetric.
  \end{block}

  \begin{proof}
    For reflexivity, let \((a,b)\in\N^2\).  The equality \(a+b=b+a\)
    shows directly that \((a,b)\sim(a,b)\).

    For symmetry, suppose \((a,b)\sim(c,d)\), so \(a+d=b+c\).  Then
    \[
      c+b=b+c=a+d=d+a.
    \]
    The first and last equalities use commutativity.  Thus
    \(c+b=d+a\), which is exactly \((c,d)\sim(a,b)\).
  \end{proof}
\end{frame}

\begin{frame}{Transitivity}
  \small
  \begin{block}{Proposition}
    If \((a,b)\sim(c,d)\) and \((c,d)\sim(e,f)\), then
    \((a,b)\sim(e,f)\).
  \end{block}

  \begin{proof}
    The assumptions say
    \[
      a+d=b+c,
      \qquad
      c+f=d+e.
    \]
    Add the two equalities and rearrange:
    \[
      (a+f)+(c+d)=(b+e)+(c+d).
    \]
    Addition in \(\N\) admits cancellation, so cancelling the common summand
    \(c+d\) yields
    \[
      a+f=b+e.
    \]
    This is precisely the condition \((a,b)\sim(e,f)\).
  \end{proof}
\end{frame}

\begin{frame}{The integers}
  Reflexivity, symmetry, and transitivity show that \(\sim\) is an equivalence
  relation.

  \begin{block}{Definition}
    The set of integers is the set of equivalence classes
    \[
      \boxed{\Z=\N^2/{\sim}}.
    \]
    We write
    \[
      \cl{a}{b}
    \]
    for the class represented by \((a,b)\).
  \end{block}

  Thus the equality criterion in \(\Z\) is
  \[
    \boxed{\cl{a}{b}=\cl{c}{d}
      \quad\Longleftrightarrow\quad a+d=b+c.}
  \]
  This criterion will be used repeatedly.
\end{frame}

\begin{frame}{The universal property of a quotient}
  \footnotesize
  Let \(X\) be a set with an equivalence relation \(\sim\), let \(Y\) be a
  set, and write \(q\colon X\to X/{\sim}\), \(q(x)=[x]\), for the quotient map.

  \begin{block}{Universal property}
    A function \(f\colon X\to Y\) descends to a function
    \(\bar f\colon X/{\sim}\to Y\) satisfying
    \[
      \bar f([x])=f(x)
    \]
    if and only if it is constant on equivalence classes:
    \[
      x\sim x'\quad\Longrightarrow\quad f(x)=f(x'),
    \]
    When it exists, \(\bar f\) is unique; equivalently,
    \(f=\bar f\circ q\).
  \end{block}

  Thus functions \(X/{\sim}\to Y\) correspond exactly to functions
  \(X\to Y\) that are constant on equivalence classes.

  \begin{alertblock}{How to define a function on a quotient}
    Give a formula on representatives and prove that equivalent representatives
    have the same image.  The descended function is then uniquely determined.
  \end{alertblock}
\end{frame}

\subsection{Arithmetic on integers}

\begin{frame}{Negation and addition}
  The familiar rules for differences suggest the definitions
  \[
    -(a-b)=b-a,
    \qquad
    (a-b)+(c-d)=(a+c)-(b+d).
  \]

  \begin{block}{Definitions}
    For integer classes, set
    \[
      -\cl{a}{b}=\cl{b}{a},
      \qquad
      \cl{a}{b}+\cl{c}{d}=\cl{a+c}{b+d}.
    \]
  \end{block}

  Before using these formulas, we must prove that their values do not depend
  on the chosen representatives.
\end{frame}

\begin{frame}{Negation is well defined}
  \begin{block}{Proposition}
    If \((a,b)\sim(c,d)\), then \((b,a)\sim(d,c)\).
  \end{block}

  \begin{proof}
    The assumption is
    \[
      a+d=b+c.
    \]
    To prove \((b,a)\sim(d,c)\), we need
    \[
      b+c=a+d,
    \]
    which is the same equality read in the opposite direction.
  \end{proof}

  Hence equivalent pairs have equivalent negatives, so
  \(-\cl{a}{b}\) is independent of the representative \((a,b)\).
\end{frame}

\begin{frame}{Addition is well defined}
  \small
  Suppose
  \[
    (a,b)\sim(a',b'),
    \qquad
    (c,d)\sim(c',d').
  \]
  Thus \(a+b'=b+a'\) and \(c+d'=d+c'\).  Adding these equations and
  rearranging gives
  \[
    (a+c)+(b'+d')=(b+d)+(a'+c').
  \]
  Therefore
  \[
    (a+c,b+d)\sim(a'+c',b'+d').
  \]

  We have proved that replacing either summand by an equivalent pair does not
  alter the class of the sum.  Consequently,
  \[
    \cl{a}{b}+\cl{c}{d}=\cl{a+c}{b+d}
  \]
  is a well-defined operation on \(\Z\).
\end{frame}

\begin{frame}{Discovering multiplication}
  Formally expanding a product of differences gives
  \[
    (a-b)(c-d)=(ac+bd)-(ad+bc).
  \]
  This suggests the following definition.

  \begin{block}{Definition}
    \[
      \boxed{
      \cl{a}{b}\cl{c}{d}
        =\cl{ac+bd}{ad+bc}.}
    \]
  \end{block}

  For example,
  \[
    \cl{5}{2}\cl{1}{4}
      =\cl{5\cdot1+2\cdot4}{5\cdot4+2\cdot1}
      =\cl{13}{22},
  \]
  which represents \(-9\), as expected from \(3\cdot(-3)=-9\).
\end{frame}

\begin{frame}{Multiplication respects one representative}
  \footnotesize
  Suppose \((a,b)\sim(c,d)\), so \(a+d=b+c\), and fix \((e,f)\).
  Multiplying the given equality by \(e\) and by \(f\) gives
  \[
    ae+de=be+ce,
    \qquad
    af+df=bf+cf.
  \]
  We must compare
  \[
    (ae+bf,af+be)
    \quad\text{and}\quad
    (ce+df,cf+de).
  \]
  Their cross-sums satisfy
  \begin{align*}
    (ae+bf)+(cf+de)
      &=(ae+de)+(bf+cf)\\
      &=(be+ce)+(af+df)\\
      &=(af+be)+(ce+df).
  \end{align*}
  Hence the two product pairs are equivalent.
\end{frame}

\begin{frame}{Multiplication is well defined}
  Pair multiplication is commutative:
  \[
    (ae+bf,af+be)=(ea+fb,eb+fa),
  \]
  by commutativity in \(\N\).

  The preceding argument therefore works in either factor.  If
  \[
    (a,b)\sim(a',b'),
    \qquad
    (c,d)\sim(c',d'),
  \]
  then
  \[
    (a,b)(c,d)\sim(a',b')(c,d)
      \sim(a',b')(c',d').
  \]
  Transitivity now shows that the product class is unchanged when both
  representatives are replaced.

  Thus multiplication is a well-defined operation on \(\Z\).
\end{frame}

\begin{frame}{Zero and one}
  Define
  \[
    0=\cl{0}{0},
    \qquad
    1=\cl{1}{0}.
  \]

  The equality criterion immediately gives a useful test:
  \[
    \boxed{\cl{a}{b}=0\quad\Longleftrightarrow\quad a=b.}
  \]
  Indeed,
  \[
    \cl{a}{b}=\cl{0}{0}
    \quad\Longleftrightarrow\quad
    a+0=b+0
    \quad\Longleftrightarrow\quad
    a=b.
  \]

  In particular \(0\ne1\), since \(0\ne1\) in \(\N\).
\end{frame}

\begin{frame}{Associativity of addition}
  \small
  \begin{block}{Theorem}
    For all \(x,y,z\in\Z\), \((x+y)+z=x+(y+z)\).
  \end{block}

  \begin{proof}
    Choose representatives
    \[
      x=\cl{a}{b},\qquad y=\cl{c}{d},\qquad z=\cl{e}{f}.
    \]
    Then
    \begin{align*}
      (x+y)+z
        &=\cl{(a+c)+e}{(b+d)+f},\\
      x+(y+z)
        &=\cl{a+(c+e)}{b+(d+f)}.
    \end{align*}
    The two displayed pairs are equal by associativity of addition in \(\N\).
    Therefore their integer classes are equal.
  \end{proof}
\end{frame}

\begin{frame}{The additive identity and additive inverses}
  Let \(x=\cl{a}{b}\).

  \begin{block}{Zero is an identity}
    \[
      x+0=\cl{a+0}{b+0}=\cl{a}{b}=x.
    \]
    Commutativity of natural addition similarly gives \(0+x=x\).
  \end{block}

  \begin{block}{Every integer has an additive inverse}
    \[
      x+(-x)=\cl{a+b}{b+a}.
    \]
    Since \(a+b=b+a\), the zero criterion gives
    \[
      x+(-x)=0.
    \]
  \end{block}
\end{frame}

\begin{frame}{Commutativity and subtraction}
  If \(x=\cl{a}{b}\) and \(y=\cl{c}{d}\), then
  \[
    x+y=\cl{a+c}{b+d}
       =\cl{c+a}{d+b}
       =y+x.
  \]
  Thus addition on \(\Z\) is commutative.

  \begin{block}{Definition of subtraction}
    We may now define
    \[
      x-y=x+(-y).
    \]
    In terms of representatives,
    \[
      \cl{a}{b}-\cl{c}{d}
        =\cl{a+d}{b+c}.
    \]
  \end{block}

  The operation that motivated the construction has now become available.
\end{frame}

\begin{frame}{The multiplicative identity and commutativity}
  Let \(x=\cl{a}{b}\).  Then
  \[
    1x
      =\cl{1}{0}\cl{a}{b}
      =\cl{1a+0b}{1b+0a}
      =\cl{a}{b}
      =x.
  \]

  If \(y=\cl{c}{d}\), then
  \begin{align*}
    xy&=\cl{ac+bd}{ad+bc},\\
    yx&=\cl{ca+db}{cb+da}.
  \end{align*}
  Corresponding coordinates are equal by commutativity in \(\N\), so \(xy=yx\).
  Consequently \(x1=1x=x\) as well.
\end{frame}

\begin{frame}{Associativity of multiplication: first coordinates}
  \footnotesize
  Let
  \[
    x=\cl{a}{b},\qquad y=\cl{c}{d},\qquad z=\cl{e}{f}.
  \]
  The first coordinate of a representative of \((xy)z\) is
  \begin{align*}
    (ac+bd)e+(ad+bc)f
      &=ace+bde+adf+bcf.
  \end{align*}
  The first coordinate of a representative of \(x(yz)\) is
  \begin{align*}
    a(ce+df)+b(cf+de)
      &=ace+adf+bcf+bde.
  \end{align*}
  These are equal by associativity, commutativity, and distributivity in \(\N\).
\end{frame}

\begin{frame}{Associativity of multiplication: second coordinates}
  \footnotesize
  With the same notation, the second coordinate of \((xy)z\) is
  \[
    (ac+bd)f+(ad+bc)e
      =acf+bdf+ade+bce.
  \]
  The second coordinate of \(x(yz)\) is
  \[
    a(cf+de)+b(ce+df)
      =acf+ade+bce+bdf.
  \]
  Again the two expressions are equal after rearranging their summands.
  Therefore the representing pairs are equal, and
  \[
    \boxed{(xy)z=x(yz).}
  \]
\end{frame}

\begin{frame}{Distributivity}
  \footnotesize
  Let \(x=\cl{a}{b}\), \(y=\cl{c}{d}\), and \(z=\cl{e}{f}\).
  The first coordinate of \(x(y+z)\) is
  \[
    a(c+e)+b(d+f)=ac+ae+bd+bf,
  \]
  while that of \(xy+xz\) is
  \[
    (ac+bd)+(ae+bf).
  \]
  They are equal after rearrangement.  The second coordinates are likewise
  \[
    a(d+f)+b(c+e)=ad+af+bc+be
  \]
  and
  \[
    (ad+bc)+(af+be).
  \]
  Hence \(x(y+z)=xy+xz\).  Commutativity of multiplication gives
  \((x+y)z=xz+yz\).
\end{frame}

\begin{frame}{The algebraic structure of \(\Z\)}
  We have proved that
  \begin{itemize}
    \item addition is associative and commutative, with identity \(0\);
    \item every integer has an additive inverse;
    \item multiplication is associative and commutative, with identity \(1\);
    \item multiplication distributes over addition;
    \item \(0\ne1\).
  \end{itemize}

  Therefore
  \[
    (\Z,0,1,+,-,\cdot)
  \]
  is a \alert{commutative ring}.

  Every one of these laws ultimately follows from the arithmetic of \(\N\) and
  the equality criterion for integer classes.
\end{frame}

\begin{frame}{The rule of signs}
  The classes \(\cl{n}{0}\) and \(\cl{0}{n}=-\cl{n}{0}\) represent a
  natural number and its negative.  Applying the multiplication formula
  directly gives, for \(m,n\in\N\),
  \[
  \begin{array}{c|cc}
    \cdot & \cl{n}{0} & \cl{0}{n}\\ \hline
    \cl{m}{0} & \cl{mn}{0} & \cl{0}{mn}\\[0.4em]
    \cl{0}{m} & \cl{0}{mn} & \cl{mn}{0}
  \end{array}
  \]

  This is the familiar rule of signs:
  \[
    (+)(+)=+,\quad (+)(-)=-,\quad (-)(+)=-,\quad (-)(-)=+.
  \]

  The table illustrates the signs of these particular representatives.
\end{frame}

\begin{frame}{A natural-number lemma}
  \small
  \begin{block}{Lemma}
    If \(a,b,c,d\in\N\) satisfy
    \[
      ad+bc=ac+bd,
    \]
    then \(a=b\) or \(c=d\).
  \end{block}

  The proof follows the induction used in the natural-number development.
  Induct on \(a\), allowing \(b\) to vary.

  If \(a=0\), the equality is \(bc=bd\), so either \(b=0=a\) or
  cancellation gives \(c=d\).  If \(a\) is a successor, split into the cases
  \(b=0\) and \(b=\suc f\).  The first again gives \(c=d\).  In the second,
  expanding both sides and cancelling the common terms \(c+d\) reduces the
  equality to the induction hypothesis.
\end{frame}

\begin{frame}{No zero divisors}
  \begin{block}{Theorem}
    If \(x\ne0\) and \(y\ne0\), then \(xy\ne0\).
  \end{block}

  \begin{proof}
    Choose arbitrary representatives
    \[
      x=\cl{a}{b},
      \qquad
      y=\cl{c}{d}.
    \]
    If \(xy=0\), the multiplication formula and the equality criterion give
    \[
      ad+bc=ac+bd.
    \]
    The natural-number lemma yields \(a=b\) or \(c=d\).  In the first case
    \(x=\cl{a}{a}=0\); in the second, \(y=\cl{c}{c}=0\).  Both conclusions
    contradict the hypotheses, so \(xy\ne0\).
  \end{proof}
\end{frame}

\begin{frame}{Cancellation for multiplication}
  \begin{block}{Theorem}
    If \(x\ne0\) and \(yx=zx\), then \(y=z\).
  \end{block}

  \begin{proof}
    Subtract \(zx\) from both sides:
    \[
      yx-zx=0.
    \]
    By distributivity,
    \[
      (y-z)x=0.
    \]
    If \(y-z\ne0\), then both factors would be nonzero, contradicting the
    no-zero-divisors theorem.  Hence \(y-z=0\), which is equivalent to \(y=z\).
  \end{proof}

  The assumption \(x\ne0\) is necessary, since \(y0=z0\) for all \(y,z\).
\end{frame}

\subsection{The natural numbers inside the integers}

\begin{frame}{The natural embedding}
  \begin{block}{Definition}
    Define
    \[
      \iota\colon\N\longrightarrow\Z,
      \qquad
      \iota(n)=\cl{n}{0}.
    \]
  \end{block}

  This map sends a natural number to the integer represented by \(n-0\).
  It preserves the distinguished elements:
  \[
    \iota(0)=\cl{0}{0}=0,
    \qquad
    \iota(1)=\cl{1}{0}=1.
  \]
\end{frame}

\begin{frame}{The embedding preserves arithmetic}
  For \(a,b\in\N\),
  \begin{align*}
    \iota(a)+\iota(b)
      &=\cl{a}{0}+\cl{b}{0}\\
      &=\cl{a+b}{0}
       =\iota(a+b),
  \end{align*}
  and
  \begin{align*}
    \iota(a)\iota(b)
      &=\cl{a}{0}\cl{b}{0}\\
      &=\cl{ab+0\cdot0}{a\cdot0+0\cdot b}\\
      &=\cl{ab}{0}
       =\iota(ab).
  \end{align*}

  Thus natural addition and multiplication agree with the integer operations.
\end{frame}

\begin{frame}{The embedding is injective}
  \begin{block}{Proposition}
    If \(\iota(a)=\iota(b)\), then \(a=b\).
  \end{block}

  \begin{proof}
    The equality
    \[
      \cl{a}{0}=\cl{b}{0}
    \]
    is equivalent to
    \[
      a+0=0+b.
    \]
    Hence \(a=b\).
  \end{proof}

  We may therefore identify \(\N\) with its image in \(\Z\), obtaining the
  familiar inclusion
  \[
    \N\subseteq\Z.
  \]
\end{frame}

\subsection{The order relation}

\begin{frame}{Order as a nonnegative gap}
  \begin{block}{Definition}
    For \(x,y\in\Z\), define
    \[
      \boxed{x\le y
      \quad\Longleftrightarrow\quad
      y=x+\iota(n)\text{ for some }n\in\N.}
    \]
  \end{block}

  Thus \(x\le y\) means that the gap \(y-x\) is a natural number.

  Immediate examples are
  \[
    x\le x
    \quad\text{using }n=0,
    \qquad
    0\le1
    \quad\text{using }n=1.
  \]
\end{frame}

\begin{frame}{Reflexivity of the integer order}
  \begin{block}{Proposition}
    For every \(x\in\Z\), one has \(x\le x\).
  \end{block}

  Choose \(0\in\N\) as the witness in the definition of \(\le\).  Since
  \(\iota(0)=0\),
  \[
    x=x+0=x+\iota(0).
  \]
  Therefore \(x\le x\).
\end{frame}

\begin{frame}{Transitivity of the integer order}
  Suppose \(x\le y\) and \(y\le z\).  Unpacking the two hypotheses gives
  \(p,q\in\N\) such that
  \[
    y=x+\iota(p),
    \qquad
    z=y+\iota(q).
  \]
  Substitute these equalities and use \(p+q\) as the new witness:
  \begin{align*}
    z
      &=(x+\iota(p))+\iota(q)\\
      &=x+\bigl(\iota(p)+\iota(q)\bigr)\\
      &=x+\iota(p+q).
  \end{align*}
  Hence \(x\le z\).
\end{frame}

\begin{frame}{Antisymmetry: the two gaps}
  Suppose \(x\le y\) and \(y\le x\).  Choose \(p,q\in\N\) such that
  \[
    y=x+\iota(p),
    \qquad
    x=y+\iota(q).
  \]
  Substitute \(y=x+\iota(p)\) in the second equality.  Then
  \[
    x=(x+\iota(p))+\iota(q)=x+\iota(p+q).
  \]
  Cancelling \(x\) from the additive equality yields
  \[
    \iota(p+q)=0=\iota(0).
  \]
\end{frame}

\begin{frame}{Antisymmetry: the gaps vanish}
  Since \(\iota\) is injective,
  \[
    \iota(p+q)=\iota(0)
    \quad\Longrightarrow\quad
    p+q=0.
  \]
  A sum of natural numbers is zero only when both summands are zero, so
  \[
    p=0,\qquad q=0.
  \]
  Therefore
  \[
    y=x+\iota(p)=x+\iota(0)=x.
  \]
  We have proved that \(x\le y\) and \(y\le x\) imply \(x=y\).
\end{frame}

\begin{frame}{Totality: the first natural comparison}
  \small
  Let
  \[
    x=\cl{a}{b},
    \qquad
    y=\cl{c}{d}.
  \]
  The natural numbers \(a+d\) and \(b+c\) are comparable:
  \[
    a+d\le b+c
    \qquad\text{or}\qquad
    b+c\le a+d.
  \]

  Suppose first that \(a+d\le b+c\).  By the definition of the order on
  \(\N\), choose \(e\in\N\) such that
  \[
    b+c=(a+d)+e
  \]
  and use \(e\) as a witness for \(x\le y\).  Indeed,
  \[
    x+\iota(e)=\cl{a+e}{b}.
  \]
  The displayed natural-number equality, after reassociation and
  commutation, is
  \[
    c+b=d+(a+e).
  \]
  Thus \(\cl{c}{d}=\cl{a+e}{b}=x+\iota(e)\), so \(x\le y\).
\end{frame}

\begin{frame}{Totality: the second natural comparison}
  \small
  Now suppose that \(b+c\le a+d\).  Choose \(e\in\N\) such that
  \[
    a+d=(b+c)+e.
  \]
  This time use \(e\) as a witness for \(y\le x\).  We have
  \[
    y+\iota(e)=\cl{c+e}{d},
  \]
  while reassociating the chosen equality gives
  \[
    a+d=b+(c+e).
  \]
  Hence
  \[
    \cl{a}{b}=\cl{c+e}{d}=y+\iota(e),
  \]
  and therefore \(y\le x\).  In either case, \(x\le y\) or \(y\le x\).
\end{frame}

\begin{frame}{The integers form a linear order}
  We have proved, directly from the witness definition of \(\le\), that
  for all \(x,y,z\in\Z\):
  \begin{align*}
    &x\le x,\\
    &x\le y\text{ and }y\le z \quad\Longrightarrow\quad x\le z,\\
    &x\le y\text{ and }y\le x \quad\Longrightarrow\quad x=y,\\
    &x\le y\text{ or }y\le x.
  \end{align*}

  \begin{block}{Conclusion}
    These are exactly the defining properties of a linear order.  Thus
    \((\Z,\le)\) is a linear order.
  \end{block}
\end{frame}

\begin{frame}{The embedding preserves and reflects order}
  \small
  \begin{block}{Proposition}
    For \(a,b\in\N\),
    \[
      \boxed{\iota(a)\le\iota(b)\Longleftrightarrow a\le b.}
    \]
  \end{block}

  Suppose \(\iota(a)\le\iota(b)\).  Choose \(n\in\N\) such that
  \[
    \iota(b)=\iota(a)+\iota(n)=\iota(a+n).
  \]
  Injectivity of \(\iota\) gives \(b=a+n\), hence \(a\le b\).

  Conversely, if \(a\le b\), choose \(k\in\N\) such that \(b=a+k\).  Then
  \[
    \iota(b)=\iota(a+k)=\iota(a)+\iota(k),
  \]
  so \(k\) witnesses \(\iota(a)\le\iota(b)\).
\end{frame}

\subsection{Order and arithmetic}

\begin{frame}{Addition preserves order}
  \begin{block}{Theorem}
    If \(x\le y\), then \(x+z\le y+z\).
  \end{block}

  Choose \(n\in\N\) such that \(y=x+\iota(n)\).  Then
  \begin{align*}
    y+z
      &=(x+\iota(n))+z\\
      &=(x+z)+\iota(n),
  \end{align*}
  where associativity and commutativity of addition justify the rearrangement.
  Thus the same natural gap \(n\) proves
  \[
    x+z\le y+z.
  \]
\end{frame}

\begin{frame}{Strict positivity}
  Since the order is linear, we write
  \[
    x<y
    \quad\Longleftrightarrow\quad
    x\le y\text{ and }x\ne y.
  \]

  If \(0<x\), then \(0\le x\), so the definition of order gives
  \[
    x=0+\iota(n)=\iota(n)
  \]
  for some \(n\in\N\).  Moreover \(n\ne0\), because \(x\ne0\) and \(\iota\)
  is injective.

  Conversely, if \(n\ne0\), then
  \[
    0<\iota(n).
  \]
  Indeed \(0\le\iota(n)\) has gap \(n\), and injectivity gives
  \(\iota(n)\ne0\).

  Therefore the positive integers are exactly the images of the positive
  natural numbers.
\end{frame}

\begin{frame}{A product of positive integers is positive}
  \begin{block}{Theorem}
    If \(0<x\) and \(0<y\), then \(0<xy\).
  \end{block}

  Write
  \[
    x=\iota(n),\qquad y=\iota(m),
    \qquad n,m\ne0.
  \]
  Since the embedding preserves multiplication,
  \[
    xy=\iota(n)\iota(m)=\iota(nm).
  \]
  The natural-number product \(nm\) is nonzero.  The characterization of
  positive integers therefore gives
  \[
    0<\iota(nm)=xy.
  \]
\end{frame}

\begin{frame}{An ordered commutative ring}
  We have shown that
  \begin{itemize}
    \item \(\Z\) is a commutative ring;
    \item \(\le\) is a linear order;
    \item adding the same integer preserves order;
    \item the product of two positive integers is positive.
  \end{itemize}

  These facts make \(\Z\) an \alert{ordered commutative ring}.

  They also recover all the familiar sign behavior.  For example, translating
  \(x<y\) by \(-x\) gives
  \[
    0<y-x.
  \]
\end{frame}

\begin{frame}{The Archimedean bound}
  \begin{block}{Theorem}
    For every \(x\in\Z\), there is an \(n\in\N\) such that
    \[
      x\le\iota(n).
    \]
  \end{block}

  Let \(x=\cl{a}{b}\).  We claim that \(n=a\) works.  Indeed,
  \[
    x+\iota(b)
      =\cl{a}{b}+\cl{b}{0}
      =\cl{a+b}{b}.
  \]
  But
  \[
    \cl{a+b}{b}=\cl{a}{0}=\iota(a),
  \]
  because \((a+b)+0=b+a\).  Thus
  \[
    \iota(a)=x+\iota(b),
  \]
  and the natural gap \(b\) proves \(x\le\iota(a)\).
\end{frame}

\subsection{What we have built}

\begin{frame}{Final picture}
  Starting from \(\N\), we constructed \(\Z\) and proved that
  \begin{itemize}
    \item integers are equivalence classes of formal differences;
    \item addition, negation, and multiplication are independent of
      representatives;
    \item \(\Z\) is a commutative ring with no zero divisors;
    \item nonzero factors may be cancelled;
    \item \(\N\) embeds into \(\Z\), preserving arithmetic and order;
    \item the nonnegative-gap relation is a linear order;
    \item addition preserves order and positive products are positive;
    \item every integer is bounded above by a natural number.
  \end{itemize}
\end{frame}


\begin{frame}{An equation without solutions in \(\Z\)}
  Suppose that some \(x\in\Z\) satisfied
  \[
    2x=1.
  \]
  By totality of the integer order,
  \[
    0\le x
    \qquad\text{or}\qquad
    x\le0.
  \]

  If \(0\le x\), write \(x=\iota(n)\) with \(n\in\N\).  Since \(\iota\)
  preserves multiplication and is injective,
  \[
    \iota(2n)=2\iota(n)=2x=1=\iota(1).
  \]
  Hence
  \[
    2n=1.
  \]

  If \(x\le0\), write \(0=x+\iota(n)\) with \(n\in\N\).  Multiplying by
  \(2\) and again using injectivity gives
  \[
    0=2x+\iota(2n)=1+\iota(2n)=\iota(1+2n),
  \]
  and hence
  \[
    0=1+2n.
  \]
\end{frame}

\begin{frame}{An equation without solutions in \(\Z\) (continued)}
  We now rule out the two natural-number equations exactly as before, using
  the recursive description of addition and the successor principles.

  First,
  \[
    1+2n=2n+1=\suc{(2n)},
  \]
  so \(0=1+2n\) is impossible because zero is not a successor.

  Now suppose that \(2n=1\).  If \(n=0\), this says \(0=\suc0\), which is
  impossible.  If \(n=\suc m\), then
  \[
    2n=2(m+1)=2m+2=\suc{\suc{(2m)}}.
  \]
  Hence \(2n=1=\suc0\) would imply, by injectivity of the successor
  operation,
  \[
    \suc{(2m)}=0,
  \]
  again impossible because zero is not a successor.

  Therefore neither natural-number equation can hold, and \(2x=1\) has no
  solution in \(\Z\).
\end{frame}

\section{The rational numbers}

\subsection{Fractions as equivalence classes}

\begin{frame}{Why introduce the rational numbers?}
  In \(\Z\), the equation
  \[
    2x=1
  \]
  has no solution.  To solve it we want to introduce fractions \(a/b\), where
  \(a,b\in\Z\) and \(b\ne0\).

  We begin with a pair \((a,b)\), intended to represent \(a/b\).

  \begin{alertblock}{The central difficulty}
    A rational number has many representatives:
    \[
      \frac12=\frac24=\frac{-3}{-6}.
    \]
    Before doing arithmetic, we must say exactly when two pairs represent the
    same number.
  \end{alertblock}
\end{frame}

\begin{frame}{The pre-rationals}
  \begin{block}{Definition}
    The set of \alert{pre-rationals} is
    \[
      \Q_{\mathrm{pre}}
        =\{(a,b)\in\Z^2:b\ne0\}.
    \]
  \end{block}

  The restriction on the denominator is essential.  It ensures that products
  of denominators remain nonzero and, later, that cross multiplication admits
  cancellation.

  We shall use fraction notation for the pair when no confusion can arise:
  \[
    (a,b)\quad\longleftrightarrow\quad \frac ab.
  \]
\end{frame}

\begin{frame}{When do two fractions agree?}
  \begin{block}{Definition}
    For \((a,b),(c,d)\in\Q_{\mathrm{pre}}\), declare
    \[
      (a,b)\sim(c,d)
      \quad\Longleftrightarrow\quad
      ad=bc.
    \]
  \end{block}

  This uses only multiplication and equality in \(\Z\), but it captures the
  familiar rule
  \[
    \frac ab=\frac cd
    \quad\Longleftrightarrow\quad ad=bc.
  \]

  For example, \((1,2)\sim(-3,-6)\) because
  \(1\cdot(-6)=2\cdot(-3)\).
\end{frame}

\begin{frame}{Reflexivity and symmetry}
  \begin{block}{Proposition}
    The relation \(\sim\) is reflexive and symmetric.
  \end{block}

  \begin{proof}
    For every \((a,b)\), commutativity gives \(ab=ba\), hence
    \((a,b)\sim(a,b)\).

    If \((a,b)\sim(c,d)\), then \(ad=bc\).  Reading the equality backwards
    and commuting both products gives
    \[
      cb=da,
    \]
    which says \((c,d)\sim(a,b)\).
  \end{proof}
\end{frame}

\begin{frame}{Transitivity uses the nonzero denominator}
  \small
  Suppose
  \[
    (a,b)\sim(c,d),
    \qquad
    (c,d)\sim(e,f).
  \]
  Thus \(ad=bc\) and \(cf=de\).  Multiplying the first equality by \(f\)
  and the second by \(b\) gives
  \[
    adf=bcf=bde.
  \]
  After rearranging,
  \[
    d(af)=d(be).
  \]
  Since \(d\ne0\), cancellation in \(\Z\) yields \(af=be\).  Therefore
  \((a,b)\sim(e,f)\).

  \begin{alertblock}{Why \(b\ne0\) belongs in the definition}
    Without a nonzero middle denominator, this cancellation argument would
    fail and transitivity would not follow.
  \end{alertblock}
\end{frame}

\begin{frame}{The rational numbers}
  We have proved that \(\sim\) is an equivalence relation.

  \begin{block}{Definition}
    The rational numbers are the equivalence classes
    \[
      \boxed{\Q=\Q_{\mathrm{pre}}/{\sim}}.
    \]
    We write
    \[
      \qcl ab
    \]
    for the class represented by \((a,b)\).
  \end{block}

  Equality in \(\Q\) is therefore governed by the single criterion
  \[
    \boxed{\qcl ab=\qcl cd\quad\Longleftrightarrow\quad ad=bc.}
  \]
\end{frame}

\subsection{Arithmetic on rational numbers}

\begin{frame}{The four basic operations}
  The usual fraction formulas suggest
  \begin{align*}
    -\qcl ab&=\qcl{-a}b,\\
    \qcl ab+\qcl cd&=\qcl{ad+bc}{bd},\\
    \qcl ab\qcl cd&=\qcl{ac}{bd}.
  \end{align*}
  If \(\qcl ab\ne0\), equivalently \(a\ne0\), we define its reciprocal by
  \[
    \left(\qcl ab\right)^{-1}
      =\qcl ba.
  \]

  All displayed denominators are nonzero: this follows from the absence of
  zero divisors in \(\Z\), and from \(a\ne0\) in the reciprocal formula.  The
  next task is to prove that the answers do not depend on the chosen
  representatives.
\end{frame}

\begin{frame}{Addition is well defined}
  \footnotesize
  Suppose
  \[
    \frac ab\sim\frac{a'}{b'},
    \qquad
    \frac cd\sim\frac{c'}{d'},
  \]
  so \(ab'=ba'\) and \(cd'=dc'\).  We compare the proposed sums by cross
  multiplication:
  \begin{align*}
    (ad+bc)b'd'
      &=ab'dd'+bb'cd'\\
      &=ba'dd'+bb'dc'\\
      &=bd(a'd'+b'c').
  \end{align*}
  Hence
  \[
    \frac{ad+bc}{bd}\sim
    \frac{a'd'+b'c'}{b'd'}.
  \]

  Thus addition passes from pre-rationals to equivalence classes.
\end{frame}

\begin{frame}{Multiplication is well defined}
  \small
  Under the same hypotheses \(ab'=ba'\) and \(cd'=dc'\),
  \begin{align*}
    (ac)(b'd')
      &=(ab')(cd')\\
      &=(ba')(dc')\\
      &=(bd)(a'c').
  \end{align*}
  Consequently,
  \[
    \frac{ac}{bd}\sim\frac{a'c'}{b'd'}.
  \]

  Negation is even more direct: from \(ad=bc\) we obtain
  \[
    (-a)d=b(-c),
  \]
  so \((-a,b)\sim(-c,d)\).

  Therefore multiplication and negation are well defined on \(\Q\).
\end{frame}

\begin{frame}{The zero class detects a zero numerator}
  \begin{block}{Proposition}
    If \(b\ne0\), then
    \[
      \qcl ab=\qcl01
      \quad\Longleftrightarrow\quad a=0.
    \]
  \end{block}

  \begin{proof}
    The equality criterion says
    \[
      a\cdot1=b\cdot0,
    \]
    which is equivalent to \(a=0\).
  \end{proof}

  In particular, a nonzero rational has a nonzero numerator in every
  representative.  This is exactly what is needed to define its reciprocal.
\end{frame}

\begin{frame}{The reciprocal is well defined}
  Suppose \((a,b)\sim(c,d)\) and that their common class is nonzero.  Then
  \(ad=bc\), and the preceding proposition gives
  \[
    a\ne0
    \qquad\text{and}\qquad
    c\ne0.
  \]
  Thus \((b,a)\) and \((d,c)\) are valid pre-rationals.  Moreover,
  \[
    (b,a)\sim(d,c)
    \quad\Longleftrightarrow\quad
    bc=ad,
  \]
  which is our original equality in reverse.

  Therefore \(x^{-1}=\qcl ba\) is independent of the representative chosen
  for every nonzero \(x\in\Q\).
\end{frame}

\begin{frame}{Zero, one, and subtraction}
  We define
  \[
    0=\qcl01,
    \qquad
    1=\qcl11.
  \]
  Subtraction is addition of the negative.  Directly from the fraction
  formulas,
  \[
    \boxed{
      \qcl ab-\qcl cd=\qcl{ad-bc}{bd}.}
  \]

  Also \(0\ne1\): otherwise the equality criterion would give
  \[
    0\cdot1=1\cdot1,
  \]
  contradicting \(0\ne1\) in \(\Z\).
\end{frame}

\begin{frame}{The commutative-ring laws}
  \small
  Every ring identity can now be checked on representatives.  For example,
  \begin{align*}
    \left(\qcl ab+\qcl cd\right)+\qcl ef
      &=\qcl{(ad+bc)f+(bd)e}{bdf}\\
      &=\qcl{adf+bcf+bde}{bdf},\\[0.3em]
    \qcl ab+\left(\qcl cd+\qcl ef\right)
      &=\qcl{a(df)+b(cf+de)}{bdf}\\
      &=\qcl{adf+bcf+bde}{bdf}.
  \end{align*}

  Associativity and commutativity of both operations, the identity laws,
  additive inverses, and distributivity all reduce in this way to polynomial
  identities in \(\Z\).

  \begin{block}{Conclusion}
    The rational numbers form a nontrivial commutative ring.
  \end{block}
\end{frame}

\begin{frame}{Multiplying by the reciprocal}
  \begin{block}{Proposition}
    If \(x\in\Q\) and \(x\ne0\), then \(xx^{-1}=1\).
  \end{block}

  \begin{proof}
    Write \(x=\qcl ab\).  Since \(x\ne0\), its numerator satisfies
    \(a\ne0\).  Therefore
    \[
      xx^{-1}
        =\qcl ab\qcl ba
        =\qcl{ab}{ba}
        =\qcl11,
    \]
    where the last equality follows from
    \((ab)\cdot1=(ba)\cdot1\).
  \end{proof}
\end{frame}

\begin{frame}{The field of rational numbers}
  We have constructed:
  \begin{itemize}
    \item a nontrivial commutative ring \(\Q\);
    \item a reciprocal \(x^{-1}\) for every nonzero rational \(x\);
    \item the identity \(xx^{-1}=1\) for every nonzero rational \(x\).
  \end{itemize}

  \begin{block}{Theorem}
    With these operations, \(\Q\) is a field.
  \end{block}

  Thus for \(y\ne0\), division can be defined by
  \[
    \frac{x}{y}=xy^{-1},
  \]
  and equations such as \(2x=1\) now have the expected solution
  \(x=2^{-1}\).
\end{frame}

\subsection{The integers inside the rationals}

\begin{frame}{The integer embedding}
  \begin{block}{Definition}
    Define
    \[
      j\colon\Z\longrightarrow\Q,
      \qquad
      j(a)=\qcl a1.
    \]
  \end{block}

  The fraction formulas show that
  \begin{align*}
    j(0)&=0,& j(1)&=1,\\
    j(a+c)&=j(a)+j(c),&
    j(ac)&=j(a)j(c).
  \end{align*}

  Hence the old integer arithmetic is preserved inside the new number system.
\end{frame}

\begin{frame}{No integers are identified}
  \begin{block}{Proposition}
    The map \(j\colon\Z\to\Q\) is injective.
  \end{block}

  \begin{proof}
    If \(j(a)=j(c)\), then
    \[
      \qcl a1=\qcl c1.
    \]
    Cross multiplication gives \(a\cdot1=1\cdot c\), hence \(a=c\).
  \end{proof}

  We may therefore regard \(\Z\) as a subring of \(\Q\), although the
  construction keeps the embedding visible.
\end{frame}

\begin{frame}{Every fraction is a quotient of integers}
  \begin{block}{Proposition}
    For \(a,b\in\Z\) with \(b\ne0\),
    \[
      \boxed{\qcl ab=j(a)j(b)^{-1}.}
    \]
  \end{block}

  Indeed,
  \[
    j(a)j(b)^{-1}
      =\qcl a1\qcl1b
      =\qcl ab.
  \]

  This formula says that the quotient construction has done exactly what was
  intended: it adjoins inverses of all nonzero integers, and every rational is
  obtained from those inverses.
\end{frame}

\begin{frame}{The natural embedding is compatible}
  Let \(\iota\colon\N\to\Z\) be the embedding already constructed for the integers.
  Define
  \[
    i\colon\N\longrightarrow\Q,
    \qquad
    i(n)=\qcl{\iota(n)}1.
  \]
  Then
  \[
    i=j\circ\iota.
  \]

  Since both earlier embeddings preserve \(0,1,+,\cdot\) and are injective,
  the same is true of \(i\).  Thus the chain
  \[
    \N\lhook\joinrel\longrightarrow
    \Z\lhook\joinrel\longrightarrow
    \Q
  \]
  preserves the arithmetic constructed at every previous stage.
\end{frame}

\subsection{Nonnegativity}

\begin{frame}{Choosing a positive denominator}
  \begin{block}{Proposition}
    Every rational number has a representative \(\qcl ab\) with \(b>0\).
  \end{block}

  \begin{proof}
    Begin with any representative \(\qcl ab\), where \(b\ne0\).  The total
    order on \(\Z\) gives either \(b>0\) or \(b<0\).

    In the first case there is nothing to do.  In the second case, \(-b>0\)
    and
    \[
      \qcl ab=\qcl{-a}{-b},
    \]
    because \(a(-b)=b(-a)\).
  \end{proof}

  Positive denominators let the sign of a rational be read from its numerator.
\end{frame}

\begin{frame}{Definition of nonnegativity}
  \begin{block}{Definition}
    A rational \(x\) is \alert{nonnegative} if there are integers \(a,b\) such
    that
    \[
      a\ge0,\qquad b>0,
      \qquad x=\qcl ab.
    \]
  \end{block}

  We write this property temporarily as \(x\in\Q_{\ge0}\).  It is independent
  of any single chosen representative: it asks only for the existence of one
  representative whose numerator and denominator have the required signs.

  Clearly \(0,1\in\Q_{\ge0}\), represented by \(0/1\) and \(1/1\).
\end{frame}

\begin{frame}{Every rational has one of the two signs}
  \begin{block}{Proposition}
    For every \(x\in\Q\), either \(x\in\Q_{\ge0}\) or
    \(-x\in\Q_{\ge0}\).
  \end{block}

  \begin{proof}
    Choose \(x=\qcl ab\) with \(b>0\).  By totality in \(\Z\), either
    \(a\ge0\) or \(a<0\).

    In the first case, \(a/b\) witnesses that \(x\) is nonnegative.  In the
    second, \(-a>0\), and
    \[
      -x=\qcl{-a}b
    \]
    is nonnegative.
  \end{proof}
\end{frame}

\begin{frame}{Opposite signs meet only at zero}
  \small
  \begin{block}{Proposition}
    If both \(x\) and \(-x\) are nonnegative, then \(x=0\).
  \end{block}

  \begin{proof}
    Choose
    \[
      x=\qcl ab,
      \qquad
      -x=\qcl cd,
    \]
    with \(a,c\ge0\) and \(b,d>0\).  Since
    \(-x=\qcl{-a}b\), equality of the second pair of representatives gives
    \[
      (-a)d=bc.
    \]
    The left side is nonpositive and the right side is nonnegative, so both
    are zero.  Since \(d>0\), this forces \(a=0\), and hence \(x=0\).
  \end{proof}
\end{frame}

\begin{frame}{Nonnegativity is closed under addition}
  Suppose
  \[
    x=\qcl ab,
    \qquad
    y=\qcl cd,
  \]
  where \(a,c\ge0\) and \(b,d>0\).  Then
  \[
    x+y=\qcl{ad+bc}{bd}.
  \]

  Products of nonnegative integers are nonnegative, so
  \(ad+bc\ge0\).  Products of positive integers are positive, so \(bd>0\).

  \begin{block}{Conclusion}
    If \(x,y\in\Q_{\ge0}\), then \(x+y\in\Q_{\ge0}\).
  \end{block}
\end{frame}

\begin{frame}{Nonnegativity and multiplication}
  With the same positive-denominator representatives,
  \[
    xy=\qcl{ac}{bd}.
  \]
  Since \(ac\ge0\) and \(bd>0\), the product is nonnegative.

  If \(x\ne0\), then \(a>0\), and
  \[
    x^{-1}=\qcl ba
  \]
  is positive because \(b>0\) and \(a>0\).

  \begin{block}{Conclusion}
    Products of nonnegative rationals are nonnegative, and reciprocals of
    positive rationals are positive.
  \end{block}
\end{frame}

\subsection{The order relation}

\begin{frame}{Order as a nonnegative gap}
  The order on \(\Z\) was defined by asking whether one integer differs from
  another by a nonnegative amount.  We use the same idea here.

  \begin{block}{Definition}
    For \(x,y\in\Q\), define
    \[
      \boxed{x\le y
        \quad\Longleftrightarrow\quad
        y-x\in\Q_{\ge0}.}
    \]
  \end{block}

  In particular,
  \[
    0\le x
    \quad\Longleftrightarrow\quad
    x\in\Q_{\ge0},
  \]
  because \(x-0=x\).
\end{frame}

\begin{frame}{Comparing positive-denominator fractions}
  \small
  \begin{block}{Proposition}
    If \(b,d>0\), then
    \[
      \boxed{
      \qcl ab\le\qcl cd
      \quad\Longleftrightarrow\quad
      ad\le bc.}
    \]
  \end{block}

  Indeed,
  \[
    \qcl cd-\qcl ab=\qcl{bc-ad}{bd},
  \]
  and \(bd>0\).  Thus the gap is nonnegative exactly when
  \(bc-ad\ge0\), which is equivalent to \(ad\le bc\).

  This recovers the ordinary cross-multiplication rule without ever reversing
  an inequality, because the denominators have first been made positive.
\end{frame}

\begin{frame}{Reflexivity and transitivity}
  \begin{block}{Proposition}
    The relation \(\le\) on \(\Q\) is reflexive and transitive.
  \end{block}

  For reflexivity,
  \[
    x-x=0\in\Q_{\ge0},
  \]
  so \(x\le x\).

  If \(x\le y\) and \(y\le z\), then \(y-x\) and \(z-y\) are
  nonnegative.  Their sum is nonnegative, and
  \[
    (y-x)+(z-y)=z-x.
  \]
  Therefore \(x\le z\).
\end{frame}

\begin{frame}{Antisymmetry}
  \begin{block}{Proposition}
    If \(x\le y\) and \(y\le x\), then \(x=y\).
  \end{block}

  \begin{proof}
    The two assumptions say that
    \[
      y-x\in\Q_{\ge0},
      \qquad
      x-y\in\Q_{\ge0}.
    \]
    But \(x-y=-(y-x)\).  A rational and its negative can both be
    nonnegative only when that rational is zero.  Hence \(y-x=0\), and
    therefore \(x=y\).
  \end{proof}
\end{frame}

\begin{frame}{Totality}
  \begin{block}{Proposition}
    For all \(x,y\in\Q\), either \(x\le y\) or \(y\le x\).
  \end{block}

  \begin{proof}
    Apply the sign dichotomy to the gap \(y-x\).  Either
    \[
      y-x\in\Q_{\ge0},
    \]
    which means \(x\le y\), or
    \[
      -(y-x)=x-y\in\Q_{\ge0},
    \]
    which means \(y\le x\).
  \end{proof}

  Together with reflexivity, transitivity, and antisymmetry, this makes
  \(\Q\) a linearly ordered set.
\end{frame}

\subsection{Order and field arithmetic}

\begin{frame}{Translation preserves order}
  \begin{block}{Proposition}
    If \(x\le y\), then \(x+t\le y+t\) for every \(t\in\Q\).
  \end{block}

  \begin{proof}
    The relevant gap does not change:
    \[
      (y+t)-(x+t)=y-x.
    \]
    Since the right side is nonnegative, so is the left side.
  \end{proof}

  Consequently, adding the same rational to both sides preserves and reflects
  both weak and strict inequalities.
\end{frame}

\begin{frame}{Products of nonnegative rationals}
  \begin{block}{Proposition}
    If \(0\le x\) and \(0\le y\), then \(0\le xy\).
  \end{block}

  This is precisely closure of \(\Q_{\ge0}\) under multiplication, translated
  through the equivalence
  \[
    0\le x\quad\Longleftrightarrow\quad x\in\Q_{\ge0}.
  \]

  Combining this fact with translation invariance yields the familiar rule:
  if \(x\le y\) and \(0\le t\), then
  \[
    xt\le yt.
  \]
\end{frame}

\begin{frame}{The embeddings preserve and reflect order}
  \begin{block}{Proposition}
    For \(p,q\in\Z\),
    \[
      j(p)\le j(q)
      \quad\Longleftrightarrow\quad
      p\le q.
    \]
  \end{block}

  Both fractions have denominator \(1>0\), so the comparison criterion gives
  \[
    \qcl p1\le\qcl q1
      \quad\Longleftrightarrow\quad
      p\cdot1\le1\cdot q.
  \]

  Since \(i=j\circ\iota\), the embedding \(i\colon\N\to\Q\) likewise preserves
  and reflects \(\le\) and \(<\).
\end{frame}

\begin{frame}{Strictly positive products}
  \begin{block}{Proposition}
    If \(0<x\) and \(0<y\), then \(0<xy\).
  \end{block}

  Choose positive-denominator representatives
  \[
    x=\qcl ab,
    \qquad
    y=\qcl cd.
  \]
  Strict positivity means \(a>0\) and \(c>0\).  Hence \(ac>0\), while
  \(bd>0\), so
  \[
    xy=\qcl{ac}{bd}>0.
  \]

  Thus multiplication by a positive rational strictly preserves order.
\end{frame}

\begin{frame}{An ordered field}
  We have now proved that \(\Q\) is simultaneously
  \begin{itemize}
    \item a field;
    \item a linear order;
    \item ordered compatibly with addition;
    \item ordered compatibly with multiplication by nonnegative elements.
  \end{itemize}

  \begin{block}{Theorem}
    The constructed rational numbers form a linearly ordered field.
  \end{block}

  Algebra and order are not separate structures here: both ultimately descend
  from the arithmetic and order of \(\Z\).
\end{frame}

\begin{frame}{The Archimedean bound: reducing to integers}
  \small
  \begin{block}{Theorem}
    For every \(x\in\Q\), there exists \(n\in\N\) such that \(x\le i(n)\).
  \end{block}

  If \(x<0\), choose \(n=0\).  Otherwise choose a representation
  \[
    x=\qcl ab
  \]
  with \(a\ge0\) and \(b>0\).

  The Archimedean property of \(\Z\) supplies \(n\in\N\) with
  \[
    a\le\iota(n).
  \]
  Since \(b\) is a positive integer, \(1\le b\), and since
  \(\iota(n)\ge0\),
  \[
    \iota(n)\le b\,\iota(n).
  \]
\end{frame}

\begin{frame}{The Archimedean bound: finishing the comparison}
  Combining the two integer inequalities gives
  \[
    a\le b\,\iota(n).
  \]
  Both denominators below are positive, so the cross-multiplication criterion
  yields
  \[
    \qcl ab
      \le
    \qcl{\iota(n)}1=i(n).
  \]

  \begin{alertblock}{What the theorem rules out}
    No rational number is larger than every natural number.  The field \(\Q\)
    contains arbitrarily large numbers.
  \end{alertblock}
\end{frame}

\subsection{What we have built}

\begin{frame}{Final picture}
  \begin{block}{Algebra}
    \(\Q\) is a field, and every element has the form \(j(a)j(b)^{-1}\) with
    \(a,b\in\Z\) and \(b\ne0\).
  \end{block}

  \begin{block}{Order}
    \(\Q\) is linearly ordered; addition preserves order, and products of
    nonnegative elements are nonnegative.
  \end{block}

  \begin{block}{Compatibility with earlier number systems}
    The embeddings \(\N\hookrightarrow\Z\hookrightarrow\Q\) preserve and
    reflect arithmetic and order, and \(\Q\) satisfies the Archimedean bound.
  \end{block}
\end{frame}

\section{\texorpdfstring{Arithmetic in \(\N\)}{Arithmetic in N}}

\subsection{Euclidean division}

\begin{frame}{Euclidean division}
  \begin{block}{Theorem}
    Let \(n,d\in\N\) with \(d>0\).  There are unique \(q,r\in\N\) such that
    \[
      \boxed{n=dq+r,\qquad r<d.}
    \]
  \end{block}

  The number \(q\) is the quotient and \(r\) is the remainder.  The bound
  \(r<d\) is what singles them out among all decompositions of \(n\).

  \begin{alertblock}{Proof strategy}
    Keep the positive divisor \(d\) fixed and use induction on the dividend
    \(n\).  Passing from \(n\) to \(\suc n\) either increments the remainder
    or resets it to zero and increments the quotient.
  \end{alertblock}
\end{frame}

\begin{frame}{Euclidean division: the induction}
  For \(n=0\), take \(q=r=0\).  Then
  \[
    0=d\cdot0+0,
    \qquad
    0<d.
  \]

  Now suppose
  \[
    n=dq+r,
    \qquad
    r<d.
  \]
  The successor test for strict order gives \(\suc r\le d\).  Hence there
  are two possibilities:
  \[
    \suc r<d
    \qquad\text{or}\qquad
    \suc r=d.
  \]

  These alternatives determine the quotient and remainder for \(\suc n\).
\end{frame}

\begin{frame}{Euclidean division: the two cases}
  If \(\suc r<d\), then
  \[
    \suc n
      =\suc{(dq+r)}
      =dq+\suc r.
  \]
  Thus the quotient stays \(q\) and the new remainder is \(\suc r\).

  If \(\suc r=d\), then
  \[
    \suc n
      =dq+\suc r
      =dq+d
      =d(q+1)+0.
  \]
  Thus the quotient becomes \(q+1\) and the new remainder is \(0<d\).

  In either case \(\suc n\) has the required decomposition, so induction
  proves existence for every \(n\in\N\).
\end{frame}

\begin{frame}{Uniqueness of quotient and remainder}
  Suppose
  \[
    n=dq+r=dq'+r',
    \qquad
    r,r'<d.
  \]
  If \(q<q'\), write \(q'=q+s\) with \(s>0\).  Then
  \[
    dq'=dq+ds,
    \qquad
    d\le ds.
  \]
  Consequently,
  \[
    n=dq+r<dq+d\le dq'+r'=n,
  \]
  a contradiction.  The same argument rules out \(q'<q\), so totality gives
  \(q=q'\).

  Addition cancellation in \(dq+r=dq+r'\) now gives \(r=r'\).  This proves
  uniqueness.
\end{frame}

\subsection{Prime numbers and Euclid's lemma}

\begin{frame}{Divisibility and prime numbers}
  For \(a,b\in\N\), define
  \[
    a\mid b
    \quad\Longleftrightarrow\quad
    b=ac\ \text{ for some }c\in\N.
  \]

  \begin{block}{Definition}
    A number \(p\in\N\) is \alert{prime} if \(2\le p\) and
    \[
      d\mid p
      \quad\Longrightarrow\quad
      d=1\ \text{ or }\ d=p.
    \]
  \end{block}

  Thus \(0\) and \(1\) are not prime, while \(2,3,5,7\) are prime.  A number
  \(n\ge2\) that is not prime is \alert{composite}; it can be written
  \[
    n=ab
    \qquad\text{with}\qquad
    2\le a<n,\qquad 2\le b<n.
  \]
\end{frame}

\begin{frame}{The greatest common divisor}
  Let \(a,b\in\N\), not both zero.

  \begin{block}{Definition}
    The \alert{greatest common divisor} of \(a\) and \(b\) is the natural
    number \(g=\gcd(a,b)\) characterized by
    \begin{enumerate}
      \item \(g\mid a\) and \(g\mid b\);
      \item if \(d\mid a\) and \(d\mid b\), then \(d\mid g\).
    \end{enumerate}
  \end{block}

  Thus \(g\) is a common divisor that is divisible by every common divisor.
  Since divisibility respects the natural order for positive numbers, it is
  also the largest common divisor in the usual order.

  For example,
  \[
    \gcd(84,30)=6.
  \]
\end{frame}

\begin{frame}{The Euclidean algorithm}
  Euclidean division can be repeated.  For \(a,b>0\), begin with
  \begin{align*}
    a&=bq_0+r_0,       &0\le r_0&<b,\\
    b&=r_0q_1+r_1,     &0\le r_1&<r_0,\\
    r_0&=r_1q_2+r_2,   &0\le r_2&<r_1,
  \end{align*}
  and continue while the remainder is nonzero.

  A strictly decreasing sequence of natural numbers must terminate.  At each
  step, the two numbers on one line have exactly the same common divisors as
  the divisor and remainder on the next line.

  \begin{block}{Conclusion}
    The last nonzero remainder is \(\gcd(a,b)\), the greatest common divisor
    of \(a\) and \(b\).
  \end{block}
\end{frame}

\begin{frame}{Bézout's identity}
  Each line of the Euclidean algorithm can be rearranged in \(\Z\):
  \[
    r_{i+1}=r_{i-1}-q_i r_i.
  \]
  Starting with the last nonzero remainder and substituting backwards
  expresses it as an integer linear combination of the original inputs.

  \begin{block}{Bézout's identity}
    For \(a,b>0\), there are \(u,v\in\Z\) such that
    \[
      \boxed{\gcd(a,b)=ua+vb.}
    \]
    Here the natural numbers are viewed inside \(\Z\).
  \end{block}

  In particular, if \(\gcd(a,b)=1\), then
  \[
    ua+vb=1
  \]
  for suitable integers \(u,v\).
\end{frame}

\begin{frame}{Euclid's lemma}
  \begin{block}{Lemma}
    If \(p\) is prime and \(p\mid ab\), then
    \[
      \boxed{p\mid a\quad\text{or}\quad p\mid b.}
    \]
  \end{block}

  Suppose \(p\nmid a\).  Every common divisor of \(p\) and \(a\) divides
  the prime \(p\), so it is \(1\) or \(p\).  The second possibility is
  excluded, hence \(\gcd(p,a)=1\).

  Bézout's identity gives \(up+va=1\).  Multiplying by \(b\),
  \[
    ubp+vab=b.
  \]
  Both terms on the left are divisible by \(p\): the first visibly, and the
  second because \(p\mid ab\).  Therefore \(p\mid b\), proving the lemma.
\end{frame}

\subsection{The fundamental theorem of arithmetic}

\begin{frame}{The fundamental theorem of arithmetic}
  \begin{block}{Theorem}
    Every \(n\in\N\) with \(n\ge2\) is a product of primes:
    \[
      n=p_1p_2\cdots p_k.
    \]
    This factorization is unique up to reordering the prime factors.
  \end{block}

  Equivalently, if
  \[
    n=p_1\cdots p_k=q_1\cdots q_\ell
  \]
  and both lists are arranged in nondecreasing order, then
  \[
    k=\ell
    \qquad\text{and}\qquad
    p_i=q_i\quad\text{for every }i.
  \]

  Existence comes from induction; uniqueness comes from Euclid's lemma.
\end{frame}

\begin{frame}{Existence of prime factorization}
  We use strong induction on \(n\), which follows from the induction principle
  for \(\N\).

  If \(n\) is prime, it is already a product consisting of one prime.  If it
  is not prime, then
  \[
    n=ab
    \qquad\text{with}\qquad
    2\le a<n,\qquad 2\le b<n.
  \]
  The induction hypothesis supplies prime factorizations
  \[
    a=p_1\cdots p_k,
    \qquad
    b=q_1\cdots q_\ell.
  \]
  Multiplying them gives
  \[
    n=p_1\cdots p_kq_1\cdots q_\ell.
  \]

  Thus every \(n\ge2\) has a prime factorization.
\end{frame}

\begin{frame}{Uniqueness of prime factorization}
  Suppose
  \[
    p_1\cdots p_k=q_1\cdots q_\ell.
  \]
  The prime \(p_1\) divides the product on the right.  Repeated application of
  Euclid's lemma shows that \(p_1\mid q_j\) for some \(j\).  Since \(q_j\) is
  prime and \(p_1\ne1\), this forces
  \[
    p_1=q_j.
  \]

  Reorder the right-hand factors so that \(q_j\) comes first and cancel this
  common nonzero factor.  Repeating the argument matches and cancels every
  prime on both sides.

  Therefore the two lists contain exactly the same primes with exactly the
  same multiplicities.  This proves uniqueness up to reordering.
\end{frame}

\subsection{An application to the rational numbers}

\begin{frame}{Prime multiplicities in a square}
  For a prime \(p\) and a positive natural number \(n\), let
  \(\nu_p(n)\) be the number of times \(p\) occurs in the unique prime
  factorization of \(n\).

  Uniqueness gives
  \[
    \nu_p(mn)=\nu_p(m)+\nu_p(n).
  \]
  In particular,
  \[
    \boxed{\nu_p(n^2)=2\nu_p(n).}
  \]

  Thus every prime occurs an even number of times in the factorization of a
  square.  For the prime \(2\), this says that \(\nu_2(n^2)\) is even.
\end{frame}

\begin{frame}{The equation \(x^2=2\) in \(\Q\)}
  Suppose that \(x\in\Q\) satisfied \(x^2=2\).  Replacing \(x\) by \(-x\)
  if necessary, we may assume \(x\ge0\).  A positive-denominator
  representative then has the form
  \[
    x=\qcl{\iota(A)}{\iota(B)},
    \qquad
    A,B\in\N,\qquad B>0.
  \]

  Squaring and using the equality criterion for rational numbers gives
  \[
    A^2=2B^2
  \]
  in \(\N\).  Both sides are positive, so their prime multiplicities are
  defined.  Counting factors of \(2\) yields
  \[
    2\nu_2(A)=\nu_2(A^2)
      =\nu_2(2B^2)=1+2\nu_2(B).
  \]
  The left side is even and the right side is odd, a contradiction.  Hence
  \(x^2=2\) has no solution in \(\Q\).
\end{frame}

\section{The real numbers}

\subsection{Cauchy sequences of rationals}

\begin{frame}{Why introduce the real numbers?}
  We proved that the equation
  \[
    x^2=2
  \]
  has no solution in \(\Q\).  Nevertheless, rational numbers can approximate
  the desired positive solution with arbitrarily small error.

  We therefore seek a larger ordered field in which
  \begin{itemize}
    \item every rational number remains visible;
    \item limits of rational Cauchy sequences exist;
    \item arithmetic and order are compatible with taking limits.
  \end{itemize}

  \begin{alertblock}{Construction}
    A real number will be an equivalence class of Cauchy sequences of rational
    numbers.  Two sequences represent the same real when their difference
    tends to zero.
  \end{alertblock}
\end{frame}

\begin{frame}{Absolute value}
  \small
  In any linearly ordered field, define
  \[
    \boxed{|x|=\max\{x,-x\}.}
  \]

  We shall use the following standard properties without reproving them.
  \begin{columns}[T,onlytextwidth]
    \begin{column}{0.48\textwidth}
      \begin{block}{Order and symmetry}
        \begin{align*}
          0&\le|x|,\\
          |x|=0&\Longleftrightarrow x=0,\\
          -|x|&\le x\le|x|,\\
          |-x|&=|x|,\\
          |x-y|&=|y-x|.
        \end{align*}
      \end{block}

      For \(\varepsilon\ge0\),
      \[
        |x|\le\varepsilon
        \Longleftrightarrow
        -\varepsilon\le x\le\varepsilon.
      \]
    \end{column}
    \begin{column}{0.48\textwidth}
      \begin{block}{Arithmetic and estimates}
        \begin{align*}
          |xy|&=|x|\,|y|,\\
          |x^{-1}|&=|x|^{-1} &&(x\ne0),\\
          |x+y|&\le|x|+|y|,\\
          \bigl||x|-|y|\bigr|&\le|x-y|.
        \end{align*}
      \end{block}

      For \(\delta\ge0\),
      \[
        \delta<|x|
        \Longleftrightarrow
        x<-\delta\ \text{ or }\ \delta<x.
      \]
    \end{column}
  \end{columns}
\end{frame}

\begin{frame}{Cauchy sequences and pre-reals}
  Let \(x=(x_n)_{n\in\N}\) be a sequence in \(\Q\).

  \begin{block}{Definition}
    The sequence \(x\) is \alert{Cauchy} if
    \[
      \forall\varepsilon>0\ \exists N\ \forall p,q\ge N,
      \qquad |x_p-x_q|\le\varepsilon.
    \]
    Here \(\varepsilon\) is rational.
  \end{block}

  The condition asks the tail of the sequence to have arbitrarily small
  diameter.  It does not mention a limit, since that limit need not lie in
  \(\Q\).

  \begin{block}{Definition}
    A \alert{pre-real} is a rational Cauchy sequence.  We write
    \[
      \R_{\mathrm{pre}}
        =\{x\colon\N\to\Q:x\text{ is Cauchy}\}.
    \]
  \end{block}
\end{frame}

\begin{frame}{Every Cauchy sequence is bounded: the finite part}
  Let \(x\) be Cauchy.  Apply the definition with \(\varepsilon=1\).  There is
  \(A\in\N\) such that
  \[
    p,q\ge A
    \quad\Longrightarrow\quad
    |x_p-x_q|\le1.
  \]

  The initial set
  \[
    \{|x_0|,|x_1|,\ldots,|x_A|\}
  \]
  is finite and nonempty, so it has a maximum \(M\).  In particular,
  \[
    0\le M,
    \qquad
    |x_n|\le M\quad(0\le n\le A).
  \]

  We claim that \(B=M+1\) bounds the whole sequence.
\end{frame}

\begin{frame}{Every Cauchy sequence is bounded: the tail}
  If \(n\le A\), then
  \[
    |x_n|\le M\le M+1.
  \]
  If \(A\le n\), the triangle inequality and the Cauchy estimate give
  \[
    |x_n|
      =|x_A+(x_n-x_A)|
      \le |x_A|+|x_n-x_A|
      \le M+1.
  \]

  Thus
  \[
    \boxed{B=M+1>0,\qquad |x_n|\le B\text{ for every }n.}
  \]

  This boundedness lemma is the key extra ingredient in proving that products
  and reciprocals of suitable Cauchy sequences are again Cauchy.
\end{frame}

\begin{frame}{When do two pre-reals represent the same number?}
  \begin{block}{Definition}
    For \(x,y\in\R_{\mathrm{pre}}\), define
    \[
      x\sim y
      \quad\Longleftrightarrow\quad
      \forall\varepsilon>0\ \exists N\ \forall n\ge N,
      \qquad |x_n-y_n|\le\varepsilon.
    \]
  \end{block}

  In words, \(x\sim y\) means that \(x_n-y_n\) tends to zero.

  Reflexivity follows from \(|x_n-x_n|=0\).  Symmetry follows from
  \[
    |y_n-x_n|=|x_n-y_n|.
  \]

  Transitivity is the only part requiring an estimate.
\end{frame}

\begin{frame}{Transitivity of the equivalence relation}
  Suppose \(x\sim y\) and \(y\sim z\), and fix \(\varepsilon>0\).  Choose
  thresholds \(N,M\) such that
  \[
    |x_n-y_n|\le\frac\varepsilon2\quad(n\ge N),
    \qquad
    |y_n-z_n|\le\frac\varepsilon2\quad(n\ge M).
  \]

  For \(n\ge\max(N,M)\),
  \begin{align*}
    |x_n-z_n|
      &=|(x_n-y_n)+(y_n-z_n)|\\
      &\le |x_n-y_n|+|y_n-z_n|\\
      &\le \frac\varepsilon2+\frac\varepsilon2
       =\varepsilon.
  \end{align*}

  Hence \(x\sim z\), and \(\sim\) is an equivalence relation on the
  pre-reals.
\end{frame}

\subsection{Arithmetic on pre-reals}

\begin{frame}{Constants, negation, and addition}
  Define operations pointwise:
  \[
    0_n=0,
    \qquad
    1_n=1,
    \qquad
    (-x)_n=-x_n,
    \qquad
    (x+y)_n=x_n+y_n.
  \]
  Constant sequences are Cauchy, and negation preserves every Cauchy
  estimate because
  \[
    |(-x_p)-(-x_q)|=|x_q-x_p|=|x_p-x_q|.
  \]

  For addition, choose \(N,M\) so that the variations of \(x\) and \(y\) are
  at most \(\varepsilon/2\).  If \(p,q\ge\max(N,M)\), then
  \begin{align*}
    |(x_p+y_p)-(x_q+y_q)|
      &\le |x_p-x_q|+|y_p-y_q|\\
      &\le\frac\varepsilon2+\frac\varepsilon2=\varepsilon.
  \end{align*}
  Thus \(-x\) and \(x+y\) are pre-reals.
\end{frame}

\begin{frame}{Addition respects equivalence}
  Suppose \(x\sim x'\) and \(y\sim y'\).  Given \(\varepsilon>0\), choose a
  common threshold after which
  \[
    |x_n-x'_n|\le\frac\varepsilon2,
    \qquad
    |y_n-y'_n|\le\frac\varepsilon2.
  \]

  Beyond that threshold,
  \begin{align*}
    |(x_n+y_n)-(x'_n+y'_n)|
      &=|(x_n-x'_n)+(y_n-y'_n)|\\
      &\le |x_n-x'_n|+|y_n-y'_n|\\
      &\le\varepsilon.
  \end{align*}

  Therefore
  \[
    x+y\sim x'+y'.
  \]
  The same argument, using symmetry of absolute value, shows that
  \(x\sim x'\) implies \(-x\sim-x'\).
\end{frame}

\begin{frame}{Products of Cauchy sequences}
  Let \(|x_n|\le A\) and \(|y_n|\le B\), with \(A,B>0\).  Choose a common
  threshold such that
  \[
    |x_p-x_q|\le\frac\varepsilon{2B},
    \qquad
    |y_p-y_q|\le\frac\varepsilon{2A}.
  \]

  Use the identity
  \[
    x_py_p-x_qy_q
      =x_p(y_p-y_q)+y_q(x_p-x_q).
  \]
  Hence
  \begin{align*}
    |x_py_p-x_qy_q|
      &\le |x_p|\,|y_p-y_q|+|y_q|\,|x_p-x_q|\\
      &\le A\frac\varepsilon{2A}+B\frac\varepsilon{2B}
       =\varepsilon.
  \end{align*}

  Thus the pointwise product \((xy)_n=x_ny_n\) is Cauchy.
\end{frame}

\begin{frame}{Multiplication respects equivalence}
  Suppose \(x\sim x'\) and \(y\sim y'\).  Choose bounds
  \[
    |x_n|\le A,
    \qquad
    |y'_n|\le B,
    \qquad A,B>0.
  \]
  The comparison of the two products is organized as
  \[
    x_ny_n-x'_ny'_n
      =x_n(y_n-y'_n)+y'_n(x_n-x'_n).
  \]

  Once
  \[
    |y_n-y'_n|\le\frac\varepsilon{2A},
    \qquad
    |x_n-x'_n|\le\frac\varepsilon{2B},
  \]
  we obtain
  \[
    |x_ny_n-x'_ny'_n|
      \le A\frac\varepsilon{2A}+B\frac\varepsilon{2B}
      =\varepsilon.
  \]
  Therefore \(xy\sim x'y'\).
\end{frame}

\begin{frame}{A nonzero pre-real stays away from zero}
  The condition \(x\not\sim0\) means that for some \(\delta>0\), no tail
  satisfies \(|x_n|\le\delta\).  Thus
  \[
    \exists\delta>0\ \forall N\ \exists M\ge N,
    \qquad \delta<|x_M|.
  \]

  Apply the Cauchy property with \(\delta/2\), obtaining a threshold \(N\).
  Choose \(M\ge N\) with \(|x_M|>\delta\).  For every \(n\ge M\),
  \begin{align*}
    |x_n|
      &=|x_M-(x_M-x_n)|\\
      &\ge |x_M|-|x_M-x_n|\\
      &>\delta-\frac\delta2=\frac\delta2.
  \end{align*}

  Hence there are \(A>0\) and \(M\) such that
  \[
    \boxed{A<|x_n|\quad\text{for every }n\ge M.}
  \]
\end{frame}

\begin{frame}{Reciprocals of nonzero pre-reals are Cauchy}
  Let \(x\not\sim0\).  By the preceding lemma, choose \(A>0\) and \(N\) with
  \(A<|x_n|\) for \(n\ge N\).  Also choose \(M\) so that
  \[
    |x_p-x_q|\le\varepsilon A^2
    \qquad(p,q\ge M).
  \]

  For \(p,q\ge\max(N,M)\), both denominators are nonzero and
  \begin{align*}
    \left|\frac1{x_p}-\frac1{x_q}\right|
      &=\frac{|x_q-x_p|}{|x_p|\,|x_q|}\\
      &\le\frac{|x_q-x_p|}{A^2}\\
      &\le\frac{\varepsilon A^2}{A^2}=\varepsilon.
  \end{align*}

  Thus \((1/x_n)\) is Cauchy whenever the class of \(x\) is nonzero.
\end{frame}

\begin{frame}{The reciprocal is independent of representatives}
  Suppose \(x\sim x'\), with neither sequence equivalent to zero.  Choose
  eventual lower bounds
  \[
    A<|x_n|,
    \qquad
    A'<|x'_n|.
  \]
  Since \(x\sim x'\), eventually
  \[
    |x_n-x'_n|\le\varepsilon AA'.
  \]

  Beyond a common threshold,
  \begin{align*}
    \left|\frac1{x_n}-\frac1{x'_n}\right|
      &=\frac{|x'_n-x_n|}{|x_n|\,|x'_n|}\\
      &\le\frac{|x'_n-x_n|}{AA'}\\
      &\le\varepsilon.
  \end{align*}

  Therefore \((1/x_n)\sim(1/x'_n)\).  Reciprocal is consequently well
  defined on every nonzero equivalence class.
\end{frame}

\subsection{The quotient field}

\begin{frame}{Definition of the real numbers}
  We have proved that \(\sim\) is an equivalence relation on
  \(\R_{\mathrm{pre}}\).

  \begin{block}{Definition}
    The real numbers are the quotient
    \[
      \boxed{\R=\R_{\mathrm{pre}}/{\sim}.}
    \]
    We write \(X=[x]\) when the real number \(X\) is represented by the
    Cauchy sequence \(x\).
  \end{block}

  If \(X=[x]\) and \(Y=[y]\), the operations descend from pointwise
  operations:
  \[
    X+Y=[x+y],
    \qquad
    -X=[-x],
    \qquad
    XY=[xy].
  \]
  If \(X=[x]\ne0\), choose \(N\) with \(x_n\ne0\) for \(n\ge N\).  Its
  reciprocal is represented by any rational sequence whose tail is
  \((1/x_n)_{n\ge N}\); the finitely many earlier values do not affect its
  class.  The estimates already proved show that every formula is independent
  of representatives.
\end{frame}

\begin{frame}{The commutative-ring laws}
  Write \(X=[x]\), \(Y=[y]\), and \(Z=[z]\).  Every ring law holds pointwise
  before passing to the quotient.  For example,
  \[
    ((x+y)+z)_n=(x_n+y_n)+z_n
      =x_n+(y_n+z_n)=(x+(y+z))_n,
  \]
  so
  \[
    (X+Y)+Z=X+(Y+Z).
  \]

  The same representative calculation proves
  \begin{itemize}
    \item associativity and commutativity of addition and multiplication;
    \item the zero and one identity laws;
    \item additive inverses;
    \item distributivity.
  \end{itemize}

  Thus \(\R\) is a commutative ring.
\end{frame}

\begin{frame}{The real numbers are nontrivial}
  The zero and one of \(\R\) are represented by the constant sequences
  \[
    (0,0,0,\ldots),
    \qquad
    (1,1,1,\ldots).
  \]

  If these classes were equal, then for \(\varepsilon=1/2\) there would be a
  threshold \(N\) such that
  \[
    |0-1|\le\frac12
    \qquad(n\ge N).
  \]
  But \(|0-1|=1\), a contradiction.  Therefore
  \[
    \boxed{0\ne1\text{ in }\R.}
  \]

  Hence the quotient contains at least two distinct elements.
\end{frame}

\begin{frame}{Multiplying by a reciprocal}
  Let \(X=[x]\ne0\).  Then \(x\not\sim0\), so there are \(\delta>0\) and
  \(N\) such that
  \[
    \delta<|x_n|
    \qquad(n\ge N).
  \]
  In particular, \(x_n\ne0\) for \(n\ge N\).  On this tail,
  \[
    x_n\frac1{x_n}=1.
  \]

  Hence the product sequence differs from the constant sequence \(1\) only
  in finitely many positions.  Their difference is eventually zero, so they
  are equivalent:
  \[
    \boxed{XX^{-1}=1\qquad(X\ne0).}
  \]

  Therefore the constructed real numbers form a field.
\end{frame}

\begin{frame}{Embedding the rational numbers}
  A rational number determines a constant Cauchy sequence.  Define
  \[
    k\colon\Q\longrightarrow\R,
    \qquad
    k(r)=[(r,r,r,\ldots)].
  \]
  Pointwise arithmetic gives
  \begin{align*}
    k(0)&=0, & k(1)&=1, & k(-r)&=-k(r),\\
    k(r+s)&=k(r)+k(s),
      & k(rs)&=k(r)k(s).
  \end{align*}

  Thus \(k\) is a field homomorphism.  It remains to check that distinct
  rationals do not become equivalent as constant sequences.
\end{frame}

\begin{frame}{The rational embedding is injective}
  Suppose \(k(r)=k(s)\).  Then the two constant sequences are equivalent:
  \[
    \forall\varepsilon>0\ \exists N\ \forall n\ge N,
    \qquad |r-s|\le\varepsilon.
  \]

  If \(r\ne s\), then \(|r-s|>0\).  Choose
  \[
    \varepsilon=\frac{|r-s|}{2}>0.
  \]
  Equivalence would give
  \[
    |r-s|\le\frac{|r-s|}{2},
  \]
  which is impossible.  Hence \(r=s\), so \(k\) is injective.

  We may therefore regard \(\Q\) as a subfield of \(\R\).
\end{frame}

\subsection{Positivity and order}

\begin{frame}{A positive pre-real}
  \begin{block}{Definition}
    A pre-real \(x\) is \alert{positive}, written \(\operatorname{Pos}(x)\), if
    \[
      \exists\delta>0\ \exists N\ \forall n\ge N,
      \qquad \delta\le x_n.
    \]
  \end{block}

  A positive pre-real is not equivalent to zero.  Indeed, if \(x\sim0\),
  fix the positive rational \(\delta/2\).  At a common large index we would have
  \[
    \delta\le x_n
    \qquad\text{and}\qquad
    |x_n|\le\frac\delta2,
  \]
  which is impossible.
\end{frame}

\begin{frame}{Why term-by-term sign conditions do not work}
  \small
  One might try to declare
  \[
    x>0\iff x_n>0\text{ for every }n,
    \qquad
    x\ge0\iff x_n\ge0\text{ for every }n.
  \]
  Neither condition gives a well-defined property of an equivalence class.

  \begin{enumerate}
    \item \emph{Changing finitely many terms changes the answer.}
      The equivalent pre-reals
      \[
        (1,1,1,\ldots)\sim(-1,1,1,\ldots)
      \]
      do not both have all terms positive.  Likewise,
      \((0,0,0,\ldots)\sim(-1,0,0,\ldots)\), but only the first has all
      terms nonnegative.

    \item \emph{Even requiring the sign only eventually is not enough.}
      Set \(u_n=1/(n+1)\).  Then every \(u_n>0\), but \(u\sim0\); the naive
      rule would therefore make the zero class positive.  Also \(-u\sim0\),
      although every term of \(-u\) is negative, so termwise nonnegativity
      would depend on the chosen representative of zero.
  \end{enumerate}

  A fixed eventual gap separates positive classes from zero.  Nonnegative
  will mean: equivalent to zero or positive.
\end{frame}

\begin{frame}{Positivity is independent of representatives}
  Suppose \(x\sim x'\) and \(x_n\ge\delta>0\) for \(n\ge N\).  Choose \(M\)
  such that
  \[
    |x_n-x'_n|\le\frac\delta2
    \qquad(n\ge M).
  \]

  For \(n\ge\max(N,M)\),
  \begin{align*}
    x'_n
      &=x_n-(x_n-x'_n)\\
      &\ge\delta-|x_n-x'_n|\\
      &\ge\delta-\frac\delta2=\frac\delta2.
  \end{align*}

  Hence \(\operatorname{Pos}(x')\).  By symmetry, positivity depends only on
  the equivalence class represented by the sequence.
\end{frame}

\begin{frame}{Positive pre-reals are closed under arithmetic}
  Suppose eventually
  \[
    x_n\ge A>0,
    \qquad
    y_n\ge B>0.
  \]
  Beyond the maximum of the two thresholds,
  \[
    x_n+y_n\ge A+B>0,
  \]
  so \(x+y\) is positive.

  Since all four quantities are nonnegative,
  \[
    x_ny_n\ge AB>0,
  \]
  so \(xy\) is positive as well.

  \begin{block}{Conclusion}
    \[
      \operatorname{Pos}(x),\operatorname{Pos}(y)
      \Longrightarrow
      \operatorname{Pos}(x+y),\operatorname{Pos}(xy).
    \]
  \end{block}
\end{frame}

\begin{frame}{Nonnegative pre-reals}
  \begin{block}{Definition}
    A pre-real \(x\) is \alert{nonnegative} if
    \[
      \operatorname{Nonneg}(x)
      \quad\Longleftrightarrow\quad
      \operatorname{Pos}(x)\ \text{ or }\ x\sim0.
    \]
  \end{block}

  This property is invariant under equivalence because both alternatives are.
  It includes every eventually pointwise nonnegative pre-real.  Indeed, if
  \(x_n\ge0\) for every \(n\ge N\), then:
  \begin{itemize}
    \item if \(x\sim0\), it is nonnegative by definition;
    \item otherwise, \(|x_n|>\delta\) eventually for some \(\delta>0\), and
      the pointwise sign forces \(x_n>\delta\).
  \end{itemize}
  Thus \(x\) is positive in the second case.
\end{frame}

\begin{frame}{Opposite nonnegative signs meet only at zero}
  Suppose both \(x\) and \(-x\) are nonnegative.  If \(x\not\sim0\), then
  also \(-x\not\sim0\), so both must be positive.  There would be
  \(\delta,\delta'>0\) and a common large index \(n\) with
  \[
    x_n\ge\delta,
    \qquad
    -x_n\ge\delta'.
  \]
  Adding yields
  \[
    0\ge\delta+\delta'>0,
  \]
  a contradiction.  Therefore
  \[
    \boxed{\operatorname{Nonneg}(x),\operatorname{Nonneg}(-x)
      \Longrightarrow x\sim0.}
  \]

  This is the antisymmetry mechanism for the real order.
\end{frame}

\begin{frame}{If a pre-real is not nonnegative}
  Assume \(x\) is not nonnegative.  Then \(x\not\sim0\) and \(x\) is not
  positive.

  From \(x\not\sim0\), choose \(\delta>0\) and \(N\) such that
  \[
    |x_n|>\delta\qquad(n\ge N).
  \]
  From the Cauchy property, choose \(M\) such that
  \[
    |x_p-x_q|\le\frac\delta2
    \qquad(p,q\ge M).
  \]

  Since \(x\) is not positive, the proposed gap \(\delta\) fails after every
  threshold.  Hence there is \(R\ge\max(N,M)\) with \(x_R<\delta\).  Together
  with \(|x_R|>\delta\), this forces
  \[
    x_R<-\delta,
    \qquad\text{so}\qquad
    -x_R>\delta.
  \]
\end{frame}

\begin{frame}{The opposite pre-real is positive}
  Continue with the index \(R\) for which \(-x_R>\delta\).  If
  \(n\ge\max(R,N,M)\), the Cauchy estimate gives
  \[
    x_R-x_n\ge-\frac\delta2.
  \]
  Therefore
  \begin{align*}
    -x_n
      &=(x_R-x_n)+(-x_R)\\
      &>-\frac\delta2+\delta
       =\frac\delta2.
  \end{align*}

  Thus \(-x\) is positive, and in particular nonnegative.  We have proved
  the sign alternative
  \[
    \boxed{\operatorname{Nonneg}(x)
      \quad\text{or}\quad
      \operatorname{Nonneg}(-x).}
  \]
  This is the key step toward totality of the order.
\end{frame}

\begin{frame}{Definition of the real order}
  Nonnegativity passes to equivalence classes.  For \(X,Y\in\R\), define
  \[
    \boxed{X\le Y
      \quad\Longleftrightarrow\quad
      Y-X\text{ is nonnegative}.}
  \]

  Reflexivity follows because \(X-X=0\).  For transitivity, if \(Y-X\) and
  \(Z-Y\) are nonnegative, their sum is nonnegative, and
  \[
    (Z-Y)+(Y-X)=Z-X.
  \]

  For antisymmetry, if \(X\le Y\) and \(Y\le X\), then both \(Y-X\) and
  \(-(Y-X)=X-Y\) are nonnegative.  The preceding lemma gives
  \[
    Y-X=0,
  \]
  hence \(X=Y\).
\end{frame}

\begin{frame}{Totality and compatibility with arithmetic}
  Apply the sign alternative to \(Y-X\).  If \(Y-X\) is nonnegative, then
  \(X\le Y\).  Otherwise \(-(Y-X)=X-Y\) is nonnegative, so \(Y\le X\).
  Thus the order is total.

  Translation preserves order because
  \[
    (Y+T)-(X+T)=Y-X.
  \]
  Products of nonnegative classes are nonnegative by the corresponding
  pre-real result.  Products of positive classes are positive.

  \begin{block}{Theorem}
    The constructed real numbers form a linearly ordered field.
  \end{block}
\end{frame}

\begin{frame}{Eventual inequalities are sufficient}
  The eventual non-strict sign condition is not necessary, but it is a useful
  sufficient condition:
  \[
    \boxed{x_n\ge0\text{ eventually}\quad\Longrightarrow\quad [x]\ge0.}
  \]
  Indeed, if \(x\sim0\), then \([x]=0\).  Otherwise, \(|x_n|>\delta\)
  eventually for some \(\delta>0\); combined with \(x_n\ge0\), this gives
  \(x_n>\delta\) eventually, so \(x\) is positive.

  Consequently, if \(X=[x]\) and \(Y=[y]\), then
  \[
    \boxed{x_n\le y_n\text{ eventually}\quad\Longrightarrow\quad X\le Y,}
  \]
  because \(y_n-x_n\ge0\) eventually.

  The converses fail.  For example, \(z_n=-1/(n+1)\) satisfies \(z\sim0\),
  so \([z]=0\ge0\), although no term of \(z\) is nonnegative.  Likewise,
  if \(x_n=0\) and \(y_n=z_n\), then \([x]=[y]\), but \(x_n\le y_n\)
  never holds.
\end{frame}

\begin{frame}{The rational embedding preserves order}
  For \(r,s\in\Q\), the difference \(k(s)-k(r)\) is represented by the
  constant sequence \(s-r\).

  If \(r<s\), this sequence is positive with the fixed gap \(s-r\).  If
  \(r=s\), it represents zero.  Hence
  \[
    r\le s\quad\Longrightarrow\quad k(r)\le k(s).
  \]

  Conversely, if \(s-r<0\), the constant sequence \(s-r\) is neither
  positive nor equivalent to zero, so it cannot be nonnegative.  Therefore
  \[
    \boxed{k(r)\le k(s)\quad\Longleftrightarrow\quad r\le s.}
  \]
  Injectivity then also gives
  \(k(r)<k(s)\Longleftrightarrow r<s\).
\end{frame}

\subsection{Rational density and convergence}

\begin{frame}{A rational approximation to any real}
  Let \(X=[x]\in\R\) and let \(\varepsilon\in\Q\) be positive.  Apply the
  Cauchy property of \(x\) with \(\varepsilon/2\).  Choose \(N\) such that
  \[
    |x_n-x_N|\le\frac\varepsilon2
    \qquad(n\ge N),
  \]
  and set \(r=x_N\).

  The two eventual inequalities
  \[
    -\frac\varepsilon2\le x_n-r\le\frac\varepsilon2
  \]
  say, through the definition of the real order, that
  \[
    -k(\varepsilon)<X-k(r)<k(\varepsilon).
  \]
  Equivalently,
  \[
    \boxed{|X-k(r)|<k(\varepsilon).}
  \]
\end{frame}

\begin{frame}{A positive rational below a positive real}
  Suppose \(X=[x]>0\).  Positivity of its representative gives
  \(\delta>0\) and \(N\) such that
  \[
    x_n\ge\delta
    \qquad(n\ge N).
  \]
  Set \(r=\delta/2\).  Then \(r>0\), and eventually
  \[
    x_n-r\ge\frac\delta2>0.
  \]
  Hence
  \[
    0<k(r)<X.
  \]

  Combining this with the preceding rational-error approximation proves the
  full density statement: for every \(E>0\) in \(\R\), choose
  \(0<k(\varepsilon)<E\), then find \(r\in\Q\) with
  \[
    \boxed{|X-k(r)|<E.}
  \]
\end{frame}

\begin{frame}{The Archimedean bound in \(\R\)}
  Let \(X=[x]\).  Boundedness of the pre-real \(x\) gives a positive rational
  \(B\) such that
  \[
    |x_m|\le B\qquad\text{for every }m.
  \]
  The Archimedean property of \(\Q\) gives \(n\in\N\) with \(B\le n+1\).
  Therefore
  \[
    x_m\le|x_m|\le B\le n+1
  \]
  for every \(m\).  Passing to classes gives
  \[
    \boxed{X\le k(n+1).}
  \]

  Thus the real construction inherits an Archimedean upper bound directly
  from the boundedness of rational Cauchy sequences.
\end{frame}

\begin{frame}{Arbitrarily small reciprocal bounds}
  Let \(E>0\).  Apply the Archimedean bound to \(|E^{-1}|\).  For some
  \(n\in\N\),
  \[
    |E^{-1}|\le k(n+1).
  \]
  Since \(E>0\), its reciprocal is positive, so the left side equals
  \(E^{-1}\).  Both sides are positive; reversing the inequality under
  reciprocals yields
  \[
    \boxed{k\!\left(\frac1{n+1}\right)\le E.}
  \]

  This is the form of the Archimedean property used in both Cauchy estimates
  below: the rational numbers \(1/(n+1)\) eventually fit below every positive
  real number.
\end{frame}

\begin{frame}{Convergence and Cauchy sequences in \(\R\)}
  For a sequence \(F\colon\N\to\R\), define
  \[
    F_n\longrightarrow X
  \]
  to mean
  \[
    \forall E>0\ \exists N\ \forall n\ge N,
    \qquad |F_n-X|\le E.
  \]

  The sequence is \alert{Cauchy} if
  \[
    \forall E>0\ \exists N\ \forall p,q\ge N,
    \qquad |F_p-F_q|\le E.
  \]

  The real numbers are \alert{complete} if every real Cauchy sequence
  converges to a real number.  The construction was designed so that a
  rational Cauchy sequence converges to the class it represents.
\end{frame}

\begin{frame}{From a pre-real to a sequence of real numbers}
  \small
  Let \(x=(x_n)_{n\in\N}\) be a pre-real.  Applying the rational embedding
  to each term gives a sequence of real numbers
  \[
    k(x_0),\ k(x_1),\ k(x_2),\ldots
  \]
  Each \(k(x_n)\) is represented by the constant row with value \(x_n\):

  \[
    \begin{array}{c|ccccc}
      \text{real number}
        & \multicolumn{5}{c}{\text{a rational representative}}\\
      \hline
      k(x_0) & x_0 & x_0 & x_0 & x_0 & \cdots\\
      k(x_1) & x_1 & x_1 & x_1 & x_1 & \cdots\\
      k(x_2) & x_2 & x_2 & x_2 & x_2 & \cdots\\
      \vdots & \vdots&\vdots&\vdots&\vdots&\\
      k(x_n) & x_n & x_n & x_n & x_n & \cdots\\
      \hline
      X=[x]             & x_0 & x_1 & x_2 & x_3 & \cdots
    \end{array}
  \]

  Fix a positive rational \(\varepsilon\).  Since \(x\) is Cauchy, there is
  \(N\) such that, whenever \(n,m\ge N\),
  \[
    |x_n-x_m|\le\varepsilon.
  \]
  Thus the \(k(x_n)\)-row and the \(X\)-row differ by at most \(\varepsilon\)
  at index \(m\) whenever \(n,m\ge N\).
\end{frame}

\begin{frame}{A representative sequence converges to its class}
  Let \(x\in\R_{\mathrm{pre}}\) and set \(X=[x]\).  We claim
  \[
    k(x_n)\longrightarrow X.
  \]

  Let us first fix a positive rational \(\varepsilon\).  The Cauchy property
  supplies \(N\) such that for \(n,m\ge N\),
  \[
    |x_n-x_m|\le\varepsilon.
  \]
  For fixed \(n\ge N\), this tail estimate is exactly the quotient-order
  statement
  \[
    |k(x_n)-X|\le k(\varepsilon).
  \]

  Now fix a positive real \(E\), and choose a positive rational
  \(\varepsilon\) with \(0<k(\varepsilon)<E\).  The same threshold then gives
  \[
    |k(x_n)-X|<E.
  \]
\end{frame}

\subsection{Completeness}

\begin{frame}{Approximating a real sequence term by term}
  Let \(F\colon\N\to\R\).  Rational density gives, for every \(n\), a rational
  \(x_n\) satisfying
  \[
    \boxed{|F_n-k(x_n)|
      <k\!\left(\frac1{n+1}\right).}
  \]

  Choosing one such rational for each \(n\) defines a sequence
  \[
    x\colon\N\longrightarrow\Q.
  \]

  To use \(x\) as a pre-real, we must prove that it is a rational Cauchy
  sequence whenever \(F\) is a real Cauchy sequence.
\end{frame}

\begin{frame}{A real sequence is a table of rational sequences}
  \small
  For every \(n\in\N\), choose a rational Cauchy sequence
  \(f^{(n)}=(f^{(n)}_m)_{m\in\N}\in\R_{\mathrm{pre}}\), with
  \(F_n=[(f^{(n)}_m)_{m\in\N}]\).
  The superscript \(n\) labels the row, while the subscript \(m\) labels a
  rational term within that row.  Rational density lets us select one
  rational \(x_n\) from sufficiently far along the \(n\)-th row.

  \[
    \begin{array}{c|cccc|c}
      \text{real term}
        & \multicolumn{4}{c|}{\text{one rational representative}}
        & \text{selected rational}\\
      \hline
      F_0=[(f^{(0)}_m)_{m\in\N}]
        & f^{(0)}_0 & f^{(0)}_1 & f^{(0)}_2 & \cdots
        & x_0=f^{(0)}_{m_0}\\
      F_1=[(f^{(1)}_m)_{m\in\N}]
        & f^{(1)}_0 & f^{(1)}_1 & f^{(1)}_2 & \cdots
        & x_1=f^{(1)}_{m_1}\\
      F_2=[(f^{(2)}_m)_{m\in\N}]
        & f^{(2)}_0 & f^{(2)}_1 & f^{(2)}_2 & \cdots
        & x_2=f^{(2)}_{m_2}\\
      \vdots & \vdots & \vdots & \vdots & & \vdots\\
      F_n=[(f^{(n)}_m)_{m\in\N}]
        & f^{(n)}_0 & f^{(n)}_1 & f^{(n)}_2 & \cdots
        & x_n=f^{(n)}_{m_n}
    \end{array}
  \]

  The selected entries form the rational sequence \(x=(x_n)\), with
  \[
    |F_n-k(x_n)|<k\!\left(\frac1{n+1}\right).
  \]
  The next argument proves that, when \(F\) is Cauchy, \(x\) is a pre-real
  and the limit of \(F\) is the class \(X=[x]\).
\end{frame}

\begin{frame}{The rational approximation is Cauchy}
  Assume \(F\) is Cauchy, and fix a positive rational \(\varepsilon\).  Choose
  \(N\) such that
  \[
    |F_p-F_q|\le\frac{k(\varepsilon)}3
    \qquad(p,q\ge N),
  \]
  and choose \(M\) such that
  \[
    k\!\left(\frac1{M+1}\right)
      \le\frac{k(\varepsilon)}3.
  \]

  If \(p,q\ge\max(M,N)\), monotonicity of reciprocals and the approximation
  bounds give
  \begin{align*}
    |k(x_p)-k(x_q)|
      &\le |k(x_p)-F_p|+|F_p-F_q|+|F_q-k(x_q)|\\
      &\le\frac{k(\varepsilon)}3
        +\frac{k(\varepsilon)}3
        +\frac{k(\varepsilon)}3
       =k(\varepsilon).
  \end{align*}
  Since \(k\) reflects order, \(|x_p-x_q|\le\varepsilon\).  Thus \(x\) is a
  rational Cauchy sequence.
\end{frame}

\begin{frame}{Every real Cauchy sequence converges}
  The rational approximation \(x\) is a pre-real, so let
  \[
    X=[x]\in\R.
  \]
  We already know that \(k(x_n)\to X\).

  Given \(E>0\), choose \(N\) so that
  \[
    |k(x_n)-X|\le\frac E2
    \qquad(n\ge N),
  \]
  and choose \(M\) so that \(k(1/(M+1))\le E/2\).  For
  \(n\ge\max(N,M)\),
  \begin{align*}
    |F_n-X|
      &\le |F_n-k(x_n)|+|k(x_n)-X|\\
      &\le\frac E2+\frac E2=E.
  \end{align*}

  \begin{block}{Completeness theorem}
    Every Cauchy sequence in \(\R\) converges to an element of \(\R\).
  \end{block}
\end{frame}

\subsection{\texorpdfstring{A solution of \(x^2=2\)}{A solution of x squared equals 2}}

\begin{frame}{Newton's rational approximations}
  Define a rational sequence \(a=(a_n)\) recursively by
  \[
    a_0=1,
    \qquad
    a_{n+1}=\frac12\left(a_n+\frac2{a_n}\right).
  \]

  The sequence is well defined and positive.  Indeed, \(a_0>0\), and if
  \(a_n>0\), then both \(a_n\) and \(2/a_n\) are positive, so their average
  is positive.

  The first few terms are
  \[
    1,\qquad
    \frac32,\qquad
    \frac{17}{12},\qquad
    \frac{577}{408},\quad\ldots
  \]
  They are rational, but they are designed so that their squares approach
  \(2\).
\end{frame}

\begin{frame}{The exact error formula}
  Define the rational error sequence \(e=(e_n)\) by
  \[
    e_n=a_n^2-2.
  \]
  Direct calculation gives
  \begin{align*}
    e_{n+1}
      &=\frac14\left(a_n+\frac2{a_n}\right)^2-2\\
      &=\frac{a_n^2+4+4/a_n^2-8}{4}\\
      &=\boxed{\frac{(a_n^2-2)^2}{4a_n^2}}
       =\frac{e_n^2}{4a_n^2}.
  \end{align*}

  Therefore \(e_{n+1}\ge0\), so
  \[
    a_n^2\ge2\qquad(n\ge1).
  \]
  Also \(a_1=3/2\), hence \(e_1=1/4\).
\end{frame}

\begin{frame}{The errors tend to zero}
  For \(n\ge1\), we have \(a_n^2\ge2\), so
  \[
    0\le e_{n+1}
      =\frac{e_n^2}{4a_n^2}
      \le\frac{e_n^2}{8}.
  \]
  Starting from \(e_1=1/4\), induction gives \(e_n\le1/4\).  Consequently,
  \[
    e_{n+1}\le\frac{e_n}{32},
  \]
  and another induction yields
  \[
    \boxed{0\le e_n\le\frac1{4\cdot32^{n-1}}\qquad(n\ge1).}
  \]

  These upper bounds tend to zero.  Concretely, \(32^m\ge m+1\) by
  induction, and the Archimedean property lets us choose \(m\) with
  \(m+1\ge1/(4\varepsilon)\).  Then
  \(1/(4\cdot32^m)\le\varepsilon\).
\end{frame}

\begin{frame}{The Newton sequence is Cauchy}
  \small
  Since \(a_n>0\) and \(a_n^2\ge2\), we have \(a_n\ge1\) for \(n\ge1\).
  Therefore
  \[
    0\le a_n-a_{n+1}
      =\frac{a_n^2-2}{2a_n}
      =\frac{e_n}{2a_n}
      \le\frac{e_n}{2}.
  \]

  If \(q>p\ge1\), telescoping and the geometric error estimate give
  \begin{align*}
    |a_p-a_q|
      &=\sum_{j=p}^{q-1}(a_j-a_{j+1})\\
      &\le\frac12\sum_{j=p}^{q-1}e_j\\
      &\le\frac{e_p}{2}\sum_{i=0}^{q-p-1}\frac1{32^i}
       <\frac{16}{31}e_p<e_p.
  \end{align*}
  Since \(e_p\to0\), this proves that \((a_n)\) is Cauchy.
\end{frame}

\begin{frame}{A real number whose square is \(2\)}
  The rational sequence \(a=(a_n)\) is a pre-real, so define
  \[
    \alpha=[a]\in\R.
  \]
  Since \(a_n\ge1\), the sequence has a fixed positive gap and \(\alpha>0\).

  Pointwise multiplication gives
  \[
    \alpha^2-2=[e].
  \]
  Given \(\varepsilon>0\), the error estimate supplies \(N\) such that for
  \(n\ge N\),
  \[
    |e_n|=e_n\le\varepsilon.
  \]
  Thus \(e\sim0\), and therefore
  \[
    \boxed{\alpha^2=2.}
  \]

  No rational has this property, so \(\alpha\) is a genuinely new number in
  the Cauchy completion of \(\Q\).
\end{frame}

\subsection{What we have built}

\begin{frame}{Final picture}
  We proved that
  \begin{itemize}
    \item pre-reals are rational Cauchy sequences, modulo differences tending
      to zero;
    \item boundedness controls products, and a nonzero class stays uniformly
      away from zero, which controls reciprocals;
    \item the quotient \(\R\) is a linearly ordered field containing an
      order-preserving copy of \(\Q\);
    \item rational points are dense in \(\R\);
    \item every real Cauchy sequence converges;
    \item the Newton Cauchy sequence defines a positive
      \(\alpha\in\R\) satisfying \(\alpha^2=2\).
  \end{itemize}
\end{frame}

\section{\texorpdfstring{The \(n\)-adic numbers}{The n-adic numbers}}

\subsection{Compatible finite approximations}

\begin{frame}{Why introduce the \(n\)-adic numbers?}
  Fix a natural number \(n\ge2\).  Instead of measuring an integer by its
  size, record its remainders modulo higher and higher powers of \(n\):
  \[
    \cdots,\qquad a\bmod n^3,\qquad a\bmod n^2,
    \qquad a\bmod n.
  \]
  Each remainder contains more information than the one before it, and the
  records agree when reduced to a smaller modulus.

  \begin{alertblock}{Construction}
    An \(n\)-adic number will be a sequence of residues
    \[
      x_k\in\Zmod{n^k}\qquad(k\in\N)
    \]
    whose entries are compatible under reduction.  In other words, it is a
    coherent choice of an integer modulo every power of \(n\).
  \end{alertblock}

  Passing to compatible records lets the finite precision grow without bound.
\end{frame}

\begin{frame}{The tower of reductions}
  Since \(n^k\mid n^{k+1}\), reduction of representatives defines a map
  \[
    \rho_k\colon\Zmod{n^{k+1}}\longrightarrow\Zmod{n^k},
    \qquad
    \rho_k([a]_{n^{k+1}})=[a]_{n^k}.
  \]
  It respects all the ring operations:
  \begin{align*}
    \rho_k(0)&=0, & \rho_k(1)&=1,\\
    \rho_k(u+v)&=\rho_k(u)+\rho_k(v), &
    \rho_k(uv)&=\rho_k(u)\rho_k(v).
  \end{align*}
  Thus every \(\rho_k\) is a ring homomorphism, and we obtain the tower
  \[
    \cdots\longrightarrow\Zmod{n^{k+1}}
      \xrightarrow{\rho_k}\Zmod{n^k}
      \longrightarrow\cdots\longrightarrow
      \Zmod n\longrightarrow\Zmod1.
  \]
  The last ring is trivial because \(n^0=1\).
\end{frame}

\begin{frame}{Compatible sequences}
  \begin{block}{Definition}
    A sequence
    \[
      x=(x_k)_{k\in\N},
      \qquad x_k\in\Zmod{n^k},
    \]
    is \alert{compatible} if
    \[
      \boxed{\rho_k(x_{k+1})=x_k\qquad\text{for every }k\in\N.}
    \]
  \end{block}

  By composing adjacent reductions, compatibility also gives
  \[
    x_j\bmod n^k=x_k
    \qquad\text{whenever }k\le j.
  \]

  So a later component determines every earlier one.  The sequence is not an
  arbitrary product of residue classes: its entries must describe the same
  number at every finite level of precision.
\end{frame}

\begin{frame}{Two examples of compatible sequences}
  Every integer \(a\) gives a compatible sequence
  \[
    x_k=[a]_{n^k}.
  \]
  Indeed, reducing \([a]_{n^{k+1}}\) modulo \(n^k\) gives \([a]_{n^k}\).

  For \(n=10\), the integer \(-1\) appears as
  \[
    0,\qquad 9,\qquad 99,\qquad 999,\qquad\ldots
  \]
  in the rings
  \[
    \Zmod1,\qquad\Zmod{10},\qquad\Zmod{10^2},
    \qquad\Zmod{10^3},\qquad\ldots
  \]
  because \(10^k-1\) reduces to \(10^{k-1}-1\) at the preceding level.

  The constant zero and constant one sequences are also compatible, where
  ``constant'' means the corresponding residue class at each level.
\end{frame}

\subsection{The ring of compatible sequences}

\begin{frame}{Compatibility survives arithmetic}
  Let \(x\) and \(y\) be compatible.  Define operations componentwise:
  \[
    (x+y)_k=x_k+y_k,
    \qquad
    (xy)_k=x_ky_k,
    \qquad
    (-x)_k=-x_k.
  \]

  Since reduction is a ring homomorphism,
  \begin{align*}
    \rho_k((x+y)_{k+1})
      &=\rho_k(x_{k+1})+\rho_k(y_{k+1})=x_k+y_k,\\
    \rho_k((xy)_{k+1})
      &=\rho_k(x_{k+1})\rho_k(y_{k+1})=x_ky_k,\\
    \rho_k((-x)_{k+1})&=-\rho_k(x_{k+1})=-x_k.
  \end{align*}
  Hence compatible sequences are closed under addition, multiplication, and
  negation; they also contain \(0\) and \(1\).
\end{frame}

\begin{frame}{Definition of the \(n\)-adic numbers}
  \begin{block}{Definition}
    The ring of \(n\)-adic numbers is
    \[
      \boxed{
      \Adic n=
      \left\{x\in\prod_{k\in\N}\Zmod{n^k}
      \;\middle|\;
      \rho_k(x_{k+1})=x_k\text{ for every }k\right\}.}
    \]
  \end{block}

  The product is a commutative ring under componentwise operations.  The
  closure properties just proved show that \(\Adic n\) is a subring.

  \begin{alertblock}{Conclusion}
    For every base \(n\), the compatible sequences form a commutative ring.
    When \(p\) is prime, this is the usual ring of \(p\)-adic integers.
  \end{alertblock}
\end{frame}

\begin{frame}{Finite-level projections}
  For each \(k\), evaluation at the \(k\)-th component defines
  \[
    \pi_k\colon\Adic n\longrightarrow\Zmod{n^k},
    \qquad
    \pi_k(x)=x_k.
  \]
  Because arithmetic is componentwise, \(\pi_k\) is a ring homomorphism:
  \begin{align*}
    \pi_k(0)&=0, & \pi_k(1)&=1,\\
    \pi_k(x+y)&=\pi_k(x)+\pi_k(y), &
    \pi_k(xy)&=\pi_k(x)\pi_k(y).
  \end{align*}

  Compatibility says that these projections commute with reduction:
  \[
    \begin{array}{ccc}
      \Adic n & \xrightarrow{\pi_{k+1}} & \Zmod{n^{k+1}}\\
      \big\Vert && \big\downarrow{\rho_k}\\
      \Adic n & \xrightarrow{\pi_k} & \Zmod{n^k}.
    \end{array}
  \]
\end{frame}

\begin{frame}{Reducing through several levels}
  Reduction is transitive.  If \(a\mid b\mid c\), then for
  \(u\in\Zmod c\),
  \[
    \bigl(u\bmod b\bigr)\bmod a=u\bmod a.
  \]
  Applying this repeatedly to a compatible sequence gives the following
  useful form of compatibility.

  \begin{block}{All-level compatibility}
    If \(x\in\Adic n\) and \(k\le j\), then
    \[
      \boxed{x_j\bmod n^k=x_k.}
    \]
  \end{block}

  In particular, if \(x_j=0\), then \(x_k=0\) for every \(k\le j\).
  Equivalently, once one component is nonzero, all later components are
  nonzero.
\end{frame}

\subsection{Prime bases and valuations}

\begin{frame}{The first nonzero component}
  Let \(p\) be prime and take a nonzero \(x\in\Adic p\).

  The component \(x_0\) lies in
  \[
    \Zmod{p^0}=\Zmod1,
  \]
  so \(x_0=0\).  On the other hand, some component must be nonzero: if every
  \(x_k\) vanished, then the whole sequence would be zero.

  Hence there is a least index
  \[
    m=\min\{k\in\N:x_k\ne0\}.
  \]
  It satisfies
  \[
    m\ge1,
    \qquad
    x_{m-1}=0,
    \qquad
    x_m\ne0.
  \]
  By all-level compatibility, \(x_j\ne0\) for every \(j\ge m\).
\end{frame}

\begin{frame}{Vanishing means divisibility}
  For \(j\ge k\), let \(\widetilde{x_j}\) be the canonical representative
  of \(x_j\), so
  \[
    0\le\widetilde{x_j}<p^j.
  \]
  All-level compatibility identifies \(x_k\) with the reduction of \(x_j\).
  Therefore
  \[
    \boxed{x_k=0
      \quad\Longleftrightarrow\quad
      p^k\mid\widetilde{x_j}.}
  \]

  In particular, if \(m\) is the first nonzero index of \(x\) and \(j\ge m\),
  then
  \[
    p^{m-1}\mid\widetilde{x_j},
    \qquad
    p^m\nmid\widetilde{x_j}.
  \]
  The representative \(\widetilde{x_j}\) is nonzero because \(x_j\ne0\).
\end{frame}

\begin{frame}{The valuation stabilizes}
  For a nonzero natural number \(A\), let \(\nu_p(A)\) be the multiplicity
  of the prime \(p\) in its factorization.  Then
  \[
    p^r\mid A
    \quad\Longleftrightarrow\quad
    r\le\nu_p(A).
  \]

  Let \(m\) be the first nonzero component of \(x\in\Adic p\).  For every
  \(j\ge m\), the previous divisibility test gives
  \[
    m-1\le\nu_p(\widetilde{x_j})
    \qquad\text{and}\qquad
    \nu_p(\widetilde{x_j})<m.
  \]
  Hence
  \[
    \boxed{\nu_p(\widetilde{x_j})=m-1\qquad(j\ge m).}
  \]

  Although the representatives change with \(j\), their \(p\)-adic
  valuation is fixed by the first nonzero component.
\end{frame}

\subsection{No zero divisors in prime base}

\begin{frame}{Two nonzero \(p\)-adic numbers}
  Let \(a,b\in\Adic p\) be nonzero, and let \(m_a,m_b\) be their first
  nonzero component indices.  Then
  \[
    m_a,m_b\ge1.
  \]
  Choose the common level
  \[
    \ell=m_a+m_b-1.
  \]
  Since \(m_a\le\ell\) and \(m_b\le\ell\), both components at level
  \(\ell\) are nonzero.  The stabilized valuation formula gives
  \[
    \nu_p(\widetilde{a_\ell})=m_a-1,
    \qquad
    \nu_p(\widetilde{b_\ell})=m_b-1.
  \]
  By uniqueness of prime factorization,
  \[
    \nu_p(\widetilde{a_\ell}\widetilde{b_\ell})
      =(m_a-1)+(m_b-1)=\ell-1.
  \]
\end{frame}

\begin{frame}{The product cannot vanish}
  Suppose for a contradiction that \(ab=0\) in \(\Adic p\).  Projecting to
  level \(\ell\) gives
  \[
    a_\ell b_\ell=0
    \qquad\text{in }\Zmod{p^\ell}.
  \]
  In terms of canonical representatives, this forces
  \[
    p^\ell\mid
      \widetilde{a_\ell}\widetilde{b_\ell}.
  \]
  Since both factors are nonzero, the prime-factorization criterion implies
  \[
    \ell
      \le\nu_p(\widetilde{a_\ell}\widetilde{b_\ell})
      =\ell-1,
  \]
  which is impossible.

  \begin{block}{No zero divisors}
    If \(p\) is prime and \(ab=0\) in \(\Adic p\), then
    \[
      \boxed{a=0\quad\text{or}\quad b=0.}
    \]
  \end{block}
\end{frame}

\begin{frame}{The prime-base ring is nontrivial}
  Consider the projection at level \(1\):
  \[
    \pi_1\colon\Adic p\longrightarrow\Zmod p.
  \]
  If \(0=1\) in \(\Adic p\), applying this ring homomorphism would give
  \[
    0=1\qquad\text{in }\Zmod p.
  \]
  But \(p\) is prime, so \(\Zmod p\) is a field and its zero and one are
  distinct.  Therefore
  \[
    \boxed{0\ne1\text{ in }\Adic p.}
  \]

  Combining nontriviality with the absence of zero divisors proves:
  \begin{alertblock}{Conclusion}
    For every prime \(p\), the ring \(\Adic p\) of \(p\)-adic integers is an
    integral domain.
  \end{alertblock}
\end{frame}

\subsection{What we have built}

\begin{frame}{Final picture}
  We proved that
  \begin{itemize}
    \item an \(n\)-adic number is a compatible sequence of residues modulo
      \(n^k\);
    \item compatible sequences form a subring of
      \(\prod_k\Zmod{n^k}\), hence a commutative ring;
    \item projection to each finite level is a ring homomorphism, and every
      later component reduces to every earlier one;
    \item for prime \(p\), the first nonzero component determines the stable
      \(p\)-adic valuation of all later representatives;
    \item those valuations add under multiplication, so two nonzero elements
      of \(\Adic p\) cannot have zero product;
    \item consequently, \(\Adic p\) is a nontrivial integral domain.
  \end{itemize}
\end{frame}

\end{document}
